\documentclass[12pt,a4paper,oneside,openany]{book}

% Компилировать через XeLaTeX.
\usepackage{fontspec}
\usepackage{polyglossia}
\setmainlanguage{russian}
\setotherlanguage{english}
\setmainfont{DejaVu Serif}[Ligatures=TeX,Scale=1.02]
\setsansfont{DejaVu Sans}[Ligatures=TeX,Scale=MatchLowercase]
\setmonofont{DejaVu Sans Mono}[Scale=MatchLowercase]

% Математика и графика
\usepackage{amsmath,amssymb,amsthm,mathtools}
\allowdisplaybreaks
\usepackage{changepage}
\usepackage{array,booktabs,longtable}
\usepackage{tikz}
\usetikzlibrary{arrows.meta,positioning,matrix,calc,decorations.pathreplacing}

% Страница и единое оформление текста
\usepackage[
  left=24mm,
  right=21mm,
  top=24mm,
  bottom=26mm,
  headsep=8mm
]{geometry}
\usepackage{indentfirst}
\usepackage{enumitem}
\usepackage{microtype}
\usepackage{xcolor}
\usepackage{setspace}
\onehalfspacing
\setlength{\parindent}{1.25em}
\setlength{\parskip}{0pt}
\emergencystretch=2em

\definecolor{logicblue}{HTML}{1F4E79}
\definecolor{logicteal}{HTML}{2F6F73}
\definecolor{logicgray}{HTML}{5F6872}
\definecolor{logicline}{HTML}{D8DEE6}
\definecolor{editorred}{HTML}{9B2C2C}
\definecolor{explainbg}{HTML}{F4F7FA}

\setlist{
  topsep=0.45\baselineskip,
  itemsep=0.2\baselineskip,
  parsep=0pt,
  leftmargin=*
}

% Оглавление
\usepackage{tocloft}
\setcounter{tocdepth}{2}
\setcounter{secnumdepth}{3}
\renewcommand{\cftchapfont}{\bfseries\color{logicblue}}
\renewcommand{\cftchappagefont}{\bfseries\color{logicblue}}
\renewcommand{\cftchapleader}{\cftdotfill{\cftdotsep}}
\setlength{\cftbeforechapskip}{0.55em}

% Заголовки
\usepackage{titlesec}
\titleformat{\chapter}[display]
  {\normalfont\sffamily\bfseries\color{logicblue}\raggedright}
  {\Large\MakeUppercase{\chaptername}\ \thechapter}
  {0.6em}
  {\Huge}
  [\vspace{0.4em}{\titlerule[0.8pt]}]
\titlespacing*{\chapter}{0pt}{-8pt}{28pt}

\titleformat{name=\chapter,numberless}[display]
  {\normalfont\sffamily\bfseries\color{logicblue}\raggedright}
  {}
  {0pt}
  {\Huge}
  [\vspace{0.4em}{\titlerule[0.8pt]}]
\titlespacing*{name=\chapter,numberless}{0pt}{-8pt}{24pt}

\titleformat{\section}
  {\normalfont\Large\sffamily\bfseries\color{logicteal}}
  {\thesection}
  {0.75em}
  {}
\titlespacing*{\section}{0pt}{1.4\baselineskip}{0.55\baselineskip}

\titleformat{\subsection}
  {\normalfont\large\sffamily\bfseries\color{logicblue}}
  {\thesubsection}
  {0.7em}
  {}
\titlespacing*{\subsection}{0pt}{1.1\baselineskip}{0.35\baselineskip}

\titleformat{\subsubsection}
  {\normalfont\normalsize\sffamily\bfseries\color{logicgray}}
  {\thesubsubsection}
  {0.65em}
  {}

% Колонтитулы
\usepackage{fancyhdr}
\pagestyle{fancy}
\fancyhf{}
\fancyhead[L]{\sffamily\small\color{logicgray}Линейная алгебра}
\fancyhead[R]{\sffamily\small\color{logicgray}\nouppercase{\leftmark}}
\fancyfoot[C]{\sffamily\small\color{logicgray}\thepage}
\renewcommand{\headrulewidth}{0.4pt}
\renewcommand{\headrule}{%
  \hbox to\headwidth{\color{logicline}\leaders\hrule height \headrulewidth\hfill}}
\setlength{\headheight}{25pt}

% Теоремные окружения: формулировки из исходного конспекта сохраняются.
\theoremstyle{definition}
\newtheorem{definition}{Определение}[chapter]
\newtheorem{example}{Пример}[chapter]

\theoremstyle{plain}
\newtheorem{theorem}{Теорема}[chapter]
\newtheorem{lemma}{Лемма}[chapter]
\newtheorem{proposition}{Утверждение}[chapter]
\newtheorem{corollary}{Следствие}[chapter]

\theoremstyle{remark}
\newtheorem{remark}{Замечание}[chapter]

% Редакторские дополнения визуально отделены от исходного слоя.
\newenvironment{explanation}
  {\begin{adjustwidth}{1.25em}{1.25em}\small
   \noindent\textcolor{logicteal}{\sffamily\bfseries Пояснение.}\ }
  {\par\end{adjustwidth}}
\newenvironment{editorialnote}
  {\begin{adjustwidth}{1.25em}{1.25em}\small
   \noindent\textcolor{editorred}{\sffamily\bfseries Редакторская пометка.}\ }
  {\par\end{adjustwidth}}
\newenvironment{visualreconstruction}
  {\begin{adjustwidth}{1.25em}{1.25em}\small
   \noindent\textcolor{logicgray}{\sffamily\bfseries Схема.}\ }
  {\par\end{adjustwidth}}

% Команды
\DeclareMathOperator{\ord}{ord}
\DeclareMathOperator{\inv}{inv}
\DeclareMathOperator{\sgn}{sgn}
\DeclareMathOperator{\Ker}{Ker}
\DeclareMathOperator{\rk}{rk}
\DeclareMathOperator{\tr}{tr}
\DeclareMathOperator{\Span}{span}
\newcommand{\N}{\mathbb{N}}
\newcommand{\Z}{\mathbb{Z}}
\newcommand{\Q}{\mathbb{Q}}
\newcommand{\R}{\mathbb{R}}
\newcommand{\C}{\mathbb{C}}
\newcommand{\F}{\mathbb{F}}
\newcommand{\eps}{\varepsilon}
\newcommand{\modelsrel}{\models}
\newcommand{\proves}{\vdash}
\newcommand{\uncertain}[1]{%
  \textcolor{editorred}{\textbf{[неуверенная реконструкция: #1]}}}

% Аргумент #1 позволяет менять ширину, по умолчанию она 1.2ex.
\newcommand{\blank}[1][1.2ex]{%
  \mbox{%
    \vrule width 0.1ex height 0.4ex depth 0pt
    \vrule width #1 height 0.1ex depth 0pt
    \vrule width 0.1ex height 0.4ex depth 0pt
  }%
}

\usepackage[hidelinks]{hyperref}

\title{\sffamily\bfseries\color{logicblue}Линейная алгебра}
\author{Ямщиков Герман M3233}
\date{}

\begin{document}
\sloppy
\frontmatter
\maketitle
\tableofcontents

\mainmatter

\chapter{Алгебраические структуры}

Алгебраическая структура состоит из множества-носителя и заданных на нём
операций. Сигнатура указывает арность каждой операции:

\[
  \mathcal A=\langle
  \underbrace{A}_{\substack{\text{носитель}\\\text{(множество)}}},
  \underbrace{f_1^{\,1},f_2^{\,2},\ldots}_{
    \substack{\text{функции}\\\text{(сигнатура)}}}
  \rangle .
\]
Верхний индекс функции указывает её арность. Иными словами, имеется
множество и заданные на нём алгебраические операции.

\section{Группоиды и группы}

Связь основных структур удобно представить схемой:

\begin{center}
\begin{tikzpicture}[
  scale=.9,
  transform shape,
  node distance=7mm and 22mm,
  every node/.style={align=center},
  arr/.style={-{Latex[length=2mm]},thick}
]
  \node (magma) {группоид (магма)};
  \node[below left=of magma] (semi) {полугруппа};
  \node[below=of semi] (monoid) {моноид};
  \node[below=of monoid] (groupL) {группа};
  \node[below right=of magma] (quasi) {квазигруппа};
  \node[below=of quasi] (loop) {лупа};
  \node[below=of loop] (groupR) {группа};
  \draw[arr] (magma) -- node[left,sloped]{\scriptsize $+$ ассоциативность} (semi);
  \draw[arr] (semi) -- node[left]{\scriptsize $+$ нейтр. эл.} (monoid);
  \draw[arr] (monoid) -- node[left]{\scriptsize $+$ симм. эл.} (groupL);
  \draw[arr] (magma) -- node[right,sloped]{\scriptsize $+$ деление} (quasi);
  \draw[arr] (quasi) -- node[right]{\scriptsize $+$ нейтр. эл.} (loop);
  \draw[arr] (loop) -- node[right]{\scriptsize $+$ ассоциативность} (groupR);
\end{tikzpicture}
\end{center}

\begin{definition}
Абелева группа --- группа с коммутативной операцией.
\end{definition}

\begin{definition}
Пусть заданы группы
\(\langle G,\cdot\rangle\) и \(\langle H,*\rangle\).
Отображение из \(G\) в \(H\) называется гомоморфизмом групп, если
\[
  f(x\cdot y)=f(x)*f(y).
\]
\end{definition}

Это условие можно изобразить диаграммой, в которой
\(v,vg,y\in G\) переходят соответственно в
\(f(v),f(vg),f(y)\in H\):

\begin{center}
\begin{tikzpicture}[>=Latex,node distance=18mm]
  \node[draw,minimum width=20mm,minimum height=25mm,label=above:$G$] (G) {};
  \node[draw,minimum width=20mm,minimum height=25mm,right=38mm of G,label=above:$H$] (H) {};
  \node at ($(G.center)+(0,7mm)$) (v) {$v$};
  \node at (G.center) (vg) {$vg$};
  \node at ($(G.center)+(0,-7mm)$) (y) {$y$};
  \node at ($(H.center)+(0,7mm)$) (fv) {$f(v)$};
  \node at (H.center) (fvg) {$f(vg)$};
  \node at ($(H.center)+(0,-7mm)$) (fy) {$f(y)$};
  \draw[->] (v) -- (fv); \draw[->] (vg) -- (fvg); \draw[->] (y) -- (fy);
\end{tikzpicture}
\end{center}

\begin{definition}
Изоморфизм --- биективный гомоморфизм. Автоморфизм --- изоморфизм
алгебраической структуры на себя. Эндоморфизм --- гомоморфизм
структуры в себя.
\end{definition}

\begin{definition}
Порядок группы (мощность) --- количество её элементов.
\end{definition}

\begin{example}
\[
  |\Z|=\infty,\qquad |\Z_n|=n,\qquad |\R^*|=\infty.
\]
\end{example}

\begin{definition}
Порядок элемента \(g\in G\) --- минимальное \(n\in\N\), для которого
\(g^n=1\). Если такого $n$ нет, порядок считается бесконечным.
\end{definition}

Для аддитивной записи это означает:
\[
  \ord(g)\text{ в }\langle\R,+\rangle,\qquad
  \ord(0)=1,\qquad \ord(v)=\infty,
\]
\[
  \underbrace{g+g+\cdots+g}_{n\text{ раз}}=ng,\qquad ng=0,
\]
и маленькая цепочка \(1\cdot0=0,\ 1\cdot g=g\).

\begin{definition}
Подсистема получается выбором подмножества носителя, замкнутого
относительно операций исходной системы.
\end{definition}


\begin{definition}
Пусть дана группа \(\langle G,*\rangle\). Её подсистема
\(\langle H,*\rangle\) называется подгруппой, если она сама является
группой относительно той же операции; при этом \(H\subseteq G\).
\end{definition}

\begin{proposition}
Подмножество \(H\subseteq G\) является подгруппой, если:
\begin{enumerate}[label=\arabic*)]
  \item \(H\) замкнута относительно групповой операции;
  \item \(v\in H\Rightarrow v^{-1}\in H\);
  \item \(e\in H\).
\end{enumerate}
\end{proposition}

\begin{proof}
Ассоциативность наследуется от $G$, а три перечисленных условия дают
замкнутую операцию, нейтральный элемент и обратный элемент для каждого
элемента $H$. Поэтому $H$ является группой относительно ограничения
операции $G$.
\end{proof}

\begin{definition}
Идемпотент --- элемент группоида, удовлетворяющий равенству \(v^2=v\).
\end{definition}

\begin{definition}
Инволюция --- элемент моноида, удовлетворяющий равенству \(v^2=e\).
\end{definition}

\section{Группа перестановок}

\(S_n\) --- группа перестановок на \(n\) точках. Например, перестановку
трёх элементов можно представить стрелочной диаграммой:

\begin{center}
\begin{tikzpicture}[>=Latex,x=10mm,y=8mm]
  \foreach \x/\lab in {0/1,1/2,2/3}{\node (t\x) at (\x,1) {$\lab$};}
  \foreach \x/\lab in {0/1,1/2,2/3}{\node (b\x) at (\x,0) {$\lab$};}
  \draw[->] (t0) -- (b1); \draw[->] (t1) -- (b2); \draw[->] (t2) -- (b0);
\end{tikzpicture}
\qquad «полная запись»
\end{center}

\[
  |S_n|=n!.
\]

\begin{definition}
Циклическая запись перестановки перечисляет последовательно элементы
каждой её орбиты. Например,

\[
\sigma=
\begin{pmatrix}
1&2&3&4&5&6&7&8&9&10\\
3&10&1&5&4&9&7&2&6&8
\end{pmatrix}
=(1\ 3)(2\ 10\ 8)(4\ 5)(6\ 9)(7).
\]
\end{definition}

Запись $(7)$ означает неподвижную точку $7\mapsto7$.

\begin{definition}
Длина цикла --- количество входящих в него элементов.
\end{definition}

\begin{lemma}
Порядок цикла длины \(k\) равен \(k\).
\end{lemma}

\begin{proof}
Каждое применение цикла переводит элемент в следующий элемент цикла.
После $k$ применений все элементы возвращаются на свои места, а раньше
этого ни один выбранный элемент не возвращается в исходную позицию.
\end{proof}

\begin{definition}
Циклы называются независимыми, если множества
перемещаемых элементов не пересекаются.
\end{definition}

\begin{theorem}
Любая перестановка раскладывается в произведение независимых циклов.
Такое разложение единственно с точностью до порядка множителей.
\end{theorem}


\begin{proof}
Для каждого элемента $i$ рассмотрим последовательность
\[
  i,\sigma(i),\sigma^2(i),\ldots.
\]
Так как множество конечно, последовательность впервые возвращается к
$i$ и образует цикл. Орбиты различных элементов либо совпадают, либо не
пересекаются, поэтому перестановка является произведением полученных
независимых циклов. Орбиты определяются самой перестановкой, что даёт
единственность с точностью до порядка циклов и циклического сдвига
записи внутри каждого из них.
\end{proof}

\begin{proposition}
Независимые циклы коммутируют.
\end{proposition}

\begin{proof}
На каждом элементе действует не более одного из двух независимых циклов.
Поэтому результат их последовательного применения не зависит от
порядка.
\end{proof}

\begin{definition}
Транспозиция --- цикл длины \(2\), \((i\ j)\).
\end{definition}

\begin{definition}
Фундаментальная транспозиция --- перестановка соседних элементов.
\end{definition}

\begin{lemma}
Любая перестановка раскладывается в произведение транспозиций.
\end{lemma}

\begin{proof}
Достаточно разложить каждый независимый цикл. Для цикла имеются,
например, два разложения:
\[
  (i_1\ i_2\ i_3\ldots i_n)
  =(i_1\ i_2)(i_2\ i_3)\cdots(i_{n-1}\ i_n),
\]
\[
  (i_1\ i_n)(i_1\ i_{n-1})\cdots(i_1\ i_3)(i_1\ i_2).
\]
\end{proof}

\begin{definition}
Для перестановки
\[
  \sigma=
  \begin{pmatrix}
  1&2&3&\cdots&n\\
  i_1&i_2&i_3&\cdots&i_n
  \end{pmatrix}
\]
пара индексов \(k<\ell\), для которой \(i_k>i_\ell\), образует
инверсию. Количество инверсий обозначается \(\inv(\sigma)\).
\end{definition}

\begin{example}
\[
  \sigma=
  \begin{pmatrix}
  1&2&3&4&5\\
  2&5&4&1&3
  \end{pmatrix}
\]
На полях инверсии подсчитаны по строкам:

\begin{center}
\begin{tabular}{c@{$\;-\;$}l}
  \(i_1=2\) & \(1\) инверсия;\\
  \(i_2=5\) & \(3\) инверсии;\\
  \(i_3=4\) & \(2\) инверсии;\\
  \(i_4=1\) & \(0\) инверсий,
\end{tabular}
\qquad поэтому \(\inv(\sigma)=6\).
\end{center}
\end{example}

\begin{definition}
Знак перестановки определяется формулой
\[
  \sgn(\sigma)=(-1)^{\inv(\sigma)}.
\]
Чётность перестановки --- чётность числа инверсий.
\end{definition}

\begin{lemma}
Фундаментальная транспозиция является нечётной перестановкой.
\end{lemma}

Следующая схема иллюстрирует фундаментальную транспозицию
\(\operatorname{swap}(1,5)\):
элементы между \(i\)-м и \(j\)-м местами сравниваются попарно; справа
учитываются инверсии с промежуточными элементами.

\begin{center}
\begin{tikzpicture}[x=11mm,y=7mm,>=Latex]
  \foreach \x/\v in {0/1,1/2,2/3,3/4,4/5}{
    \node[draw,circle,inner sep=1.2pt] (a\x) at (\x,1) {\scriptsize\v};
  }
  \foreach \x/\v in {0/5,1/2,2/3,3/4,4/1}{
    \node[draw,circle,inner sep=1.2pt] (b\x) at (\x,0) {\scriptsize\v};
  }
  \draw[->,bend left=18] (a0) to (b4);
  \draw[->,bend right=18] (a4) to (b0);
  \node[anchor=west] at (5.2,0.5) {\(\operatorname{swap}(1,5)\)};
\end{tikzpicture}
\end{center}

Под рисунком в исходной записи стоит
\(j-2i-1=2(j-i+1)+1\in2\N+1\)
\uncertain{промежуточная формула повреждена}.

\begin{proof}
При $i<j$ транспозицию $(i\ j)$ можно разложить в соседние
транспозиции:
\[
  (i\ j)
  =(i\ i+1)\cdots(j-2\ j-1)(j-1\ j)
   (j-2\ j-1)\cdots(i\ i+1).
\]
Здесь $2(j-i)-1$ множителей, то есть нечётное число. Следовательно,
фундаментальная транспозиция нечётна.
\end{proof}

\begin{editorialnote}
Именно число $2(j-i)-1$, а не повреждённая промежуточная запись,
согласуется с изображённым разложением. Вывод конспекта сохраняется:
любая перестановка раскладывается в произведение фундаментальных
транспозиций.
\end{editorialnote}



\begin{theorem}
Для любых $\sigma,\pi\in S_n$
\[
  \sgn(\sigma\pi)=\sgn(\sigma)\sgn(\pi).
\]
\end{theorem}

\begin{proof}
Умножение справа на соседнюю транспозицию меняет местами два соседних
значения перестановки. Их взаимный порядок меняется, а порядок каждой
из них относительно остальных значений сохраняет суммарную чётность.
Поэтому число инверсий меняет чётность и знак умножается на $-1$.

Разложим $\pi$ в произведение $\ell$ соседних транспозиций. Последовательно
применяя сказанное сначала к единичной перестановке, а затем к $\sigma$,
получаем
\[
  \sgn(\pi)=(-1)^\ell,\qquad
  \sgn(\sigma\pi)=(-1)^\ell\sgn(\sigma)
  =\sgn(\sigma)\sgn(\pi).
\]
\end{proof}

\begin{corollary}
$\sgn(\sigma^{-1})=\sgn(\sigma)$.
\end{corollary}

\begin{proof}
$\sigma\sigma^{-1}=\sigma^{-1}\sigma=\eps$, поэтому
$1=\sgn(\eps)=\sgn(\sigma)\sgn(\sigma^{-1})$. Оба сомножителя равны
$\pm1$, откуда следует утверждение.
\end{proof}

\begin{corollary}
Чётность цикла длины $k$ совпадает с чётностью числа $k-1$.
\end{corollary}

\begin{proof}
Цикл раскладывается в $k-1$ транспозиций:
\[
  (a_1\,a_2\,\ldots\,a_k)
  =(a_1\,a_k)(a_1\,a_{k-1})\cdots(a_1\,a_2).
\]
По мультипликативности знака его знак равен $(-1)^{k-1}$.
\end{proof}

Пусть перестановка разложена в $s$ независимых циклов длины
$k_1,\ldots,k_s$ и при этом переставляет $n=k_1+\cdots+k_s$ элементов.
Тогда
\[
  (k_1-1)+\cdots+(k_s-1)=n-s.
\]

\begin{definition}
Дефектом перестановки называется число
\[
  d(\sigma)=n-s,
\]
где $n$ --- число переставляемых элементов, а $s$ --- число независимых
циклов в её циклическом разложении. Следовательно,
$\sgn(\sigma)=(-1)^{d(\sigma)}$.
\end{definition}

\begin{theorem}
Умножение перестановки на транспозицию изменяет её чётность.
\end{theorem}

В доказательстве рассматриваются четыре возможных положения точек $i,j$
транспозиции $(i\ j)$ относительно циклов перестановки.


\begin{proof}
При умножении на $(i\ j)$ возможны четыре случая.
\begin{enumerate}[label=\arabic*)]
  \item Ни $i$, ни $j$ не входят в прежние циклы. Добавляются два
  переставляемых элемента и один цикл, поэтому дефект возрастает на один.
  \item $i$ и $j$ находятся в разных циклах. Эти циклы сливаются в один;
  число циклов уменьшается на один, поэтому дефект возрастает на один.
  \item $i$ и $j$ находятся в одном цикле. Цикл распадается на два;
  число циклов возрастает на один, поэтому дефект уменьшается на один.
  \item Только одна из точек $i,j$ входит в цикл. Добавляется один
  переставляемый элемент, число циклов не меняется, поэтому дефект
  возрастает на один.
\end{enumerate}
Во всех случаях чётность дефекта меняется.
\end{proof}

\begin{theorem}[Кэли]
Любая конечная группа порядка $n$ изоморфна некоторой подгруппе группы
$S_n$.
\end{theorem}

\begin{proof}
Пусть $G$ --- группа, $|G|=n$, и $a\in G$. Рассмотрим левое
преобразование
\[
  L_a\colon G\to G,\qquad g\mapsto ag.
\]
Если $e,g_1,\ldots,g_{n-1}$ --- все элементы $G$, то
$a,ag_1,\ldots,ag_{n-1}$ --- те же элементы, возможно, в другом порядке.
Действительно, из $ag_i=ag_j$ после умножения слева на $a^{-1}$ следует
$g_i=g_j$. Значит, $L_a$ --- биекция, то есть перестановка элементов $G$.


Для $a,b\in G$ имеем
\[
  L_{ab}(g)=(ab)g=a(bg)=L_a(L_b(g)),
\]
то есть $L_{ab}=L_a\circ L_b$, $L_e=\mathrm{id}_G$ и
$(L_a)^{-1}=L_{a^{-1}}$.

Множество
\[
  H=\{L_e,L_{g_1},\ldots,L_{g_{n-1}}\}\subseteq S_n
\]
замкнуто относительно композиции и является группой. Отображение
\[
\varphi\colon G\to H,\qquad a\mapsto L_a,
\]
биективно и сохраняет операцию, поэтому это искомый изоморфизм.
\end{proof}

\section{Классы эквивалентности и факторгруппы}

Рассмотрим аддитивную группу $\langle\Z_n,+\rangle$. Для $x,y\in\Z$
пишут
\[
  x\equiv y\pmod n
  \quad\Longleftrightarrow\quad
  n\mid(x-y).
\]
Это отношение рефлексивно, симметрично и транзитивно, следовательно,
является отношением эквивалентности. Оно разбивает $\Z$ на
непересекающиеся классы эквивалентности. Например, при $n=3$ получаются
три класса:
\[
\begin{array}{c|ccc}
 & [0]&[1]&[2]\\ \hline
\text{представители}&\ldots,-3,0,3,6,9,\ldots&
\ldots,-2,1,4,7,10,\ldots&
\ldots,-1,2,5,8,11,\ldots
\end{array}
\]

\begin{definition}
Множество классов эквивалентности называется фактор-множеством.
\end{definition}


Пусть теперь $H\leq G$. Определим
\[
  g_1\equiv g_2\pmod H
  \quad\Longleftrightarrow\quad
  g_1^{-1}g_2\in H,
\]
то есть ``частное'' элементов лежит в $H$.

Это отношение эквивалентности:
\begin{enumerate}[label=\arabic*)]
  \item $g^{-1}g=e\in H$;
  \item если $g_1^{-1}g_2\in H$, то
  $(g_1^{-1}g_2)^{-1}=g_2^{-1}g_1\in H$;
  \item если $g_1^{-1}g_2\in H$ и $g_2^{-1}g_3\in H$, то
  \[
    (g_1^{-1}g_2)(g_2^{-1}g_3)=g_1^{-1}g_3\in H.
  \]
\end{enumerate}
Следовательно, это отношение разбивает $G$ на непересекающиеся классы
эквивалентности.

\begin{definition}
Левый и правый смежные классы подгруппы $H$ имеют вид
\[
  gH=\{gh\mid h\in H\},
  \qquad
  Hg=\{hg\mid h\in H\}.
\]
\end{definition}

Смежные классы обладают следующими свойствами:
\begin{enumerate}[label=\arabic*)]
  \item левые (соответственно правые) смежные классы образуют разбиение
  группы $G$;
  \item отображение $H\to gH$, $h\mapsto gh$, является биекцией, поэтому
  $|H|=|gH|$.
\end{enumerate}


Множество левых смежных классов обозначается $G/H$.

Разбиение группы \(G\) на смежные классы можно изобразить полосами;
класс \(H\) содержит \(e\), а класс \(gH\) --- элемент \(g\):

\begin{center}
\begin{tikzpicture}[x=12mm,y=9mm]
  \draw[thick] (0,0) rectangle (5,2.4);
  \draw (1.4,0) -- (1.4,2.4);
  \draw (3.2,0) -- (3.2,2.4);
  \node at (0.7,2.65) {$H$};
  \node at (2.3,-0.3) {$gH$};
  \node at (4.55,2.65) {$G$};
  \fill (0.7,1.65) circle (1.3pt) node[left] {$e$};
  \fill (2.3,1.1) circle (1.3pt) node[right] {$g$};
  \foreach \yy in {0.35,0.65,0.95,1.25,1.55,1.85,2.15}
    \draw[->,gray] (0.8,\yy) -- (2.15,\yy);
\end{tikzpicture}
\end{center}

\begin{theorem}[Лагранжа]
Если $G$ --- конечная группа и $H\leq G$, то
\[
  |G|=|H|\,|G/H|.
\]
\end{theorem}

\begin{proof}
Левые смежные классы подгруппы $H$ попарно не пересекаются и образуют
разбиение $G$. Для каждого $g\in G$ отображение
$H\to gH$, $h\mapsto gh$, биективно, поэтому каждый класс содержит
$|H|$ элементов. Число классов равно $|G/H|$. Складывая их мощности,
получаем $|G|=|H|\,|G/H|$.
\end{proof}

\begin{definition}
Индекс подгруппы $H$ --- число её смежных классов; он обозначается
$[G:H]$.
\end{definition}

\begin{definition}
Подгруппа $H\leq G$ называется нормальной, если для любого $g\in G$
\[
  gH=Hg.
\]
Обозначение: $H\trianglelefteq G$.
\end{definition}

Следующие условия равносильны:
\begin{enumerate}[label=\arabic*)]
  \item $H\trianglelefteq G$;
  \item $gHg^{-1}=H$ для всех $g\in G$;
  \item $ghg^{-1}\in H$ для всех $g\in G$ и $h\in H$.
\end{enumerate}

\begin{definition}
Элемент $ghg^{-1}$ называется элементом, сопряжённым с $h$ при помощи
$g$. Сопряжённость является отношением эквивалентности.
\end{definition}

\begin{definition}
Факторгруппа группы $G$ по нормальной подгруппе $H$ --- множество
смежных классов $G/H$ с операцией
\[
  (aH)(bH)=(ab)H,
  \qquad
  AB=\{ab\mid a\in A,\ b\in B\}.
\]
\end{definition}

Пусть $\varphi\colon G\to H$ --- гомоморфизм. Тогда
\[
  \Im\varphi=\{\varphi(g)\mid g\in G\}\leq H,
  \qquad
  \Ker\varphi=\{g\in G\mid\varphi(g)=e\}\trianglelefteq G.
\]
Действительно, образ содержит нейтральный элемент и замкнут относительно
произведений и обратных элементов; для ядра нормальность следует из
\[
  \varphi(aga^{-1})=
  \varphi(a)\varphi(g)\varphi(a)^{-1}=e
  \quad(g\in\Ker\varphi).
\]


\begin{definition}
Для $H\trianglelefteq G$ канонической проекцией называется отображение
\[
  \pi\colon G\to G/H,
  \qquad g\mapsto gH.
\]
Оно является гомоморфизмом, причём $\Im\pi=G/H$ и $\Ker\pi=H$.
\end{definition}

\begin{theorem}[Основная теорема о гомоморфизме]
Пусть $\varphi\colon G\to H$ --- гомоморфизм. Тогда существует
изоморфизм
\[
  f\colon G/\Ker\varphi\xrightarrow{\sim}\Im\varphi,
  \qquad
  f(g\Ker\varphi)=\varphi(g),
\]
и $\varphi=f\circ\pi$, где $\pi$ --- каноническая проекция.
\end{theorem}

\begin{proof}
Положим $K=\Ker\varphi$. Для $h=\varphi(g)\in\Im\varphi$ полный
прообраз равен
\[
  \varphi^{-1}(h)=gK.
\]
В самом деле, любой $g'\in G$ записывается как $g'=gk$ с
$k=g^{-1}g'$. Тогда
\[
  \varphi(g')=\varphi(g)\varphi(k)=h
  \quad\Longleftrightarrow\quad
  \varphi(k)=e
  \quad\Longleftrightarrow\quad k\in K.
\]
Следовательно, элементы образа взаимно однозначно соответствуют смежным
классам ядра, а указанное отображение $f$ является изоморфизмом.
В частности, для конечных групп
\[
  |G|=|\Im\varphi|\,|\Ker\varphi|.
\]
\end{proof}

Теорему иллюстрирует схема: группа \(G\) разбита на классы \(gK\), \(K\)
и другие классы; все точки одного класса отображаются в один элемент
образа. Отмечены \(g\mapsto h\), \(K\mapsto e\) и
\(G\mapsto\Im\varphi\):

\begin{center}
\begin{tikzpicture}[>=Latex,x=10mm,y=8mm]
  \draw[thick] (0,1.2) rectangle (5,4);
  \draw (1.5,1.2) -- (1.5,4);
  \draw (3,1.2) -- (3,4);
  \node at (5.25,3.85) {$G$};
  \node at (0.75,4.25) {$gK$};
  \node at (2.25,4.25) {$K$};
  \fill (0.75,2.8) circle (1.2pt) node[left] {$g$};
  \fill (2.25,2.8) circle (1.2pt) node[right] {$e$};
  \draw[thick] (2.5,0) ellipse (3.6 and 0.75);
  \node at (6.35,0) {$H$};
  \node at (4.7,0.3) {$\Im\varphi$};
  \fill (0.7,0) circle (1.2pt) node[below] {$h$};
  \fill (2.25,0) circle (1.2pt) node[below] {$e$};
  \draw[->] (0.75,2.8) -- (0.7,0.15);
  \draw[->] (2.25,2.8) -- (2.25,0.15);
  \draw[->] (4.15,2.8) -- (4.1,0.15);
\end{tikzpicture}
\end{center}


\chapter{Кольца и поля}

\begin{definition}
Кольцом называется алгебраическая система $\langle R,+,\cdot\rangle$ с
двумя бинарными операциями, удовлетворяющая условиям:
\begin{enumerate}[label=\arabic*)]
  \item $\langle R,+\rangle$ --- абелева группа;
  \item для любых $x,y,z\in R$
  \[
    (x+y)z=xz+yz,
    \qquad
    x(y+z)=xy+xz.
  \]
\end{enumerate}
\end{definition}

\begin{definition}
Подмножество $L\subseteq R$ называется подкольцом, если
\begin{enumerate}[label=\arabic*)]
  \item $L$ является подгруппой аддитивной группы кольца $R$;
  \item $L$ замкнуто относительно умножения.
\end{enumerate}
Обозначение: $L\leq R$.
\end{definition}

\begin{definition}
Элемент $x\in R$ называется обратимым, если существует $x^{-1}$, для
которого $xx^{-1}=x^{-1}x=1$. Множество всех обратимых элементов кольца
$R$ обозначается $R^*$.
\end{definition}

\begin{definition}
Поле --- коммутативное ассоциативное кольцо с единицей, в котором каждый
ненулевой элемент обратим.
\end{definition}

В поле нет нетривиальных делителей нуля: если $ab=0$ и $b\ne0$, то после
умножения на $b^{-1}$ получаем $a=0$.

\begin{definition}
Подмножество $K$ поля $F$ называется подполем, если
\begin{enumerate}[label=\arabic*)]
  \item $1\in K$;
  \item $x\in K$, $x\ne0\Rightarrow x^{-1}\in K$;
  \item $K$ является подкольцом $F$.
\end{enumerate}
\end{definition}

\begin{definition}
Порядок поля --- мощность его носителя.
\end{definition}

\begin{definition}
Гомоморфизм колец $\varphi\colon R\to Q$ сохраняет обе операции:
\[
  \varphi(x+y)=\varphi(x)+\varphi(y),
  \qquad
  \varphi(xy)=\varphi(x)\varphi(y).
\]
\end{definition}


\begin{definition}
Характеристикой поля $F$ называется наименьшее $n\in\N$, для которого
\[
  \underbrace{1+\cdots+1}_{n}=0.
\]
Если такого $n$ нет, полагают $\operatorname{char}F=0$.
\end{definition}

\begin{theorem}
Если $\operatorname{char}F>0$, то $\operatorname{char}F$ --- простое
число.
\end{theorem}

\begin{proof}
Пусть $\operatorname{char}F=n$ и $n=km$, где $1<k,m<n$. Тогда
\[
  \underbrace{1+\cdots+1}_{km}
  =\left(\underbrace{1+\cdots+1}_{k}\right)
   \left(\underbrace{1+\cdots+1}_{m}\right)=0.
\]
В поле нет нетривиальных делителей нуля, поэтому один из множителей равен
нулю, что противоречит минимальности $n$.
\end{proof}

\section{Кольцо вычетов}

Группа классов вычетов
\[
  \Z_n=\{\bar0,\bar1,\bar2,\ldots,\overline{n-1}\}
  \cong\Z/n\Z
\]
становится кольцом, если определить
\[
  \bar a+\bar b=\overline{a+b},
  \qquad
  \bar a\bar b=\overline{ab}.
\]
Операции корректно определены: из $a_i\equiv b_i\pmod n$ следует
\[
  a_1+a_2\equiv b_1+b_2\pmod n,
  \qquad
  a_1a_2\equiv b_1b_2\pmod n.
\]

\begin{theorem}
$\Z_n$ является полем тогда и только тогда, когда $n$ --- простое число.
\end{theorem}

\begin{proof}
Если $n=km$ составно, где $1<k,m<n$, то
$\bar k\ne\bar0$, $\bar m\ne\bar0$, но $\bar k\bar m=\bar0$; значит,
в $\Z_n$ есть нетривиальные делители нуля, и оно не является полем.


Пусть теперь $n$ простое и $\bar a\ne\bar0$. Рассмотрим
\[
  \bar a,\ 2\bar a,\ 3\bar a,\ldots,(n-1)\bar a,\bar0.
\]
Это все элементы $\Z_n$: если $0\leq k<\ell<n$ и
$k\bar a=\ell\bar a$, то $n\mid(k-\ell)a$. Из простоты $n$ и
$n\nmid a$ следует $n\mid(k-\ell)$, что невозможно. Следовательно,
среди указанных кратных встречается $\bar1$, и у $\bar a$ существует
обратный элемент.
\end{proof}

\begin{theorem}[Бином Ньютона в поле характеристики $p$]
Если $\operatorname{char}F=p>0$, то для любых $a,b\in F$
\[
  (a+b)^p=a^p+b^p.
\]
\end{theorem}

\begin{proof}
В разложении
\[
  (a+b)^p=\sum_{k=0}^{p}\binom pk a^k b^{p-k}
\]
все промежуточные коэффициенты делятся на $p$, поскольку
\[
  \binom pk=\frac{p!}{k!(p-k)!},\qquad 0<k<p,
\]
и $p$ простое. В поле характеристики $p$ они равны нулю.
\end{proof}

\begin{corollary}[Малая теорема Ферма]
Для $a\in\Z$ и простого $p$
\[
  a^p\equiv a\pmod p.
\]
Если $a$ взаимно просто с $p$, то
\[
  a^{p-1}\equiv1\pmod p.
\]
\end{corollary}

\begin{proof}
В $\Z_p$ первое равенство следует из предыдущей теоремы (или из того,
что $\bar a$ есть сумма $a$ единиц), а второе --- умножением на
$\bar a^{-1}$.
\end{proof}


\section{Поле комплексных чисел}

На множестве $\R^2$ определим операции
\[
  (a,b)+(c,d)=(a+c,b+d),
\]
\[
  (a,b)(c,d)=(ac-bd,ad+bc).
\]
Относительно этих операций $\R^2$ является полем: нуль равен $(0,0)$,
единица --- $(1,0)$, умножение коммутативно и ассоциативно. Для
$(a,b)\ne(0,0)$
\[
  (a,b)^{-1}=\left(\frac{a}{a^2+b^2},-
  \frac{b}{a^2+b^2}\right).
\]
Действительно, из $(a,b)(x,y)=(1,0)$ получаем систему
\[
  ax-by=1,\qquad ay+bx=0,
\]
откуда следуют указанные $x,y$.

Полученное поле называется полем комплексных чисел $\C$. Поле $\R$
вкладывается в него отображением
\[
  \R\hookrightarrow\C,
  \qquad r\mapsto(r,0).
\]
Мнимая единица $i=(0,1)$ удовлетворяет $i^2=(-1,0)=-1$. Поэтому
\[
  (a,b)=a(1,0)+b(0,1)=a+bi.
\]
Это алгебраическая форма комплексного числа $z=a+bi$, где
$a=\Re z$, $b=\Im z$.

\begin{definition}
Сопряжённое число: $\bar z=a-bi$. Модуль и норма:
\[
  |z|=\sqrt{a^2+b^2},
  \qquad N(z)=|z|^2=a^2+b^2=z\bar z.
\]
\end{definition}

\begin{explanation}
Геометрически $z$ и $\bar z$ симметричны относительно вещественной оси,
а модуль $|z|$ равен длине радиус-вектора:
\end{explanation}

\begin{center}
\begin{tikzpicture}[>=Latex,scale=0.9]
  \draw[->] (-0.4,0) -- (4.2,0) node[right] {$\Re(z)$};
  \draw[->] (0,-2.2) -- (0,2.5) node[above] {$\Im(z)$};
  \draw[->,thick] (0,0) -- (2.8,1.6) node[above right] {$z$};
  \draw[->,thick] (0,0) -- (2.8,-1.6) node[below right] {$\bar z$};
  \draw[dashed] (2.8,1.6) -- (2.8,-1.6);
  \node[above,sloped] at (1.35,0.95) {$|z|$};
\end{tikzpicture}
\end{center}


Сопряжение является автоморфизмом поля $\C$:
\[
  \overline{z_1+z_2}=\bar z_1+\bar z_2,
  \qquad
  \overline{z_1z_2}=\bar z_1\bar z_2.
\]

\begin{definition}
Тригонометрическая форма комплексного числа:
\[
  z=a+bi
   =\sqrt{a^2+b^2}
    \left(\frac{a}{\sqrt{a^2+b^2}}
    +i\frac{b}{\sqrt{a^2+b^2}}\right)
   =|z|(\cos\varphi+i\sin\varphi)
   =\rho(\cos\varphi+i\sin\varphi).
\]
\end{definition}


\chapter{Кольцо многочленов}

Пусть $R$ --- кольцо. Кольцо $R[x]$ состоит из последовательностей
\[
  (a_0,a_1,a_2,\ldots),\qquad a_i\in R,
\]
в которых лишь конечное число элементов отлично от нуля. Операции
задаются формулами
\[
  (a_0,a_1,\ldots)+(b_0,b_1,\ldots)
  =(a_0+b_0,a_1+b_1,\ldots),
\]
\[
  (a_0,a_1,\ldots)(b_0,b_1,\ldots)
  =(c_0,c_1,\ldots),
  \qquad
  c_k=\sum_{i+j=k}a_i b_j=\sum_{i=0}^{k}a_i b_{k-i}.
\]
Например,
\[
  c_0=a_0b_0,
  \quad c_1=a_0b_1+a_1b_0,
  \quad c_2=a_0b_2+a_1b_1+a_2b_0.
\]

Если $R$ --- кольцо, то $R[x]$ тоже кольцо. Если $R$ ассоциативно,
коммутативно или имеет единицу, то тем же свойством обладает $R[x]$.
Если $R$ --- поле, $R[x]$ всё же, вообще говоря, не поле.

Последовательность $(a_0,a_1,a_2,\ldots)$ записывают как
\[
  a_0+a_1x+a_2x^2+\cdots.
\]
При этом
\[
  1=(1,0,0,\ldots),\quad
  x=(0,1,0,\ldots),\quad
  x^2=(0,0,1,0,\ldots).
\]

\begin{definition}
Степень ненулевого многочлена --- номер его старшего ненулевого
коэффициента. Полагают $\deg0=-\infty$.
\end{definition}

Для $f,g\in R[x]$
\[
  \deg(f+g)\leq\max\{\deg f,\deg g\},
  \qquad
  \deg(fg)\leq\deg f+\deg g.
\]


\begin{definition}
Гомоморфизм колец
$f\colon\langle R,+,\cdot\rangle\to\langle Q,+,\cdot\rangle$ сохраняет
сложение и умножение:
\[
  f(x+y)=f(x)+f(y),
  \qquad f(xy)=f(x)f(y).
\]
\end{definition}

\begin{definition}
Для $c\in R$ отображение подстановки (эвалюации)
\[
  \operatorname{ev}_c\colon R[x]\to R,
  \qquad f\mapsto f(c),
\]
является гомоморфизмом колец. Оно позволяет рассматривать многочлен как
функцию. При каждой повторной подстановке получается то же значение.
\end{definition}

\begin{editorialnote}
Формулировка верна для коммутативного кольца $R$; в некоммутативном
случае достаточно потребовать, чтобы $c$ лежал в центре $R$. Без этого
равенство $(fg)(c)=f(c)g(c)$ может нарушаться.
\end{editorialnote}

\begin{definition}
Число $c$ называется корнем многочлена $f$, если $f(c)=0$.
\end{definition}

Для многочлена с вещественными коэффициентами вместе с комплексным корнем
$c$ корнем является и $\bar c$.

\begin{definition}
Поле называется алгебраически замкнутым, если любой многочлен положительной
степени над ним имеет корень.
\end{definition}

\begin{theorem}[Основная теорема алгебры]
Поле $\C$ алгебраически замкнуто.
\end{theorem}

\begin{explanation}
Полное доказательство этой теоремы требует результатов комплексного
анализа или топологии и выходит за рамки алгебраического материала
первого семестра. В дальнейшем используется только её следствие: всякий
многочлен положительной степени над $\C$ имеет комплексный корень и,
повторяя деление на линейные множители, полностью раскладывается над
$\C$.
\end{explanation}


\chapter{Кольцо квадратных матриц}

Пусть $F$ --- поле. Отображение
\[
  A\colon I\times J\to F
\]
называется матрицей над $F$. Если $I=\{1,\ldots,m\}$ и
$J=\{1,\ldots,n\}$, матрицу записывают таблицей $A=(a_{ij})$, где
$i$ --- номер строки, $j$ --- номер столбца. Множество всех таких матриц
обозначается $M_{m\times n}(F)$; при $m=n$ пишут $M_n(F)$.

Для $A,B\in M_{m\times n}(F)$ сумма определяется покомпонентно:
\[
  C=A+B,\qquad c_{ij}=a_{ij}+b_{ij}.
\]
Если $A\in M_{m\times k}(F)$ и $B\in M_{k\times n}(F)$, то
\[
  C=AB\in M_{m\times n}(F),
  \qquad
  c_{ij}=\sum_{\ell=1}^{k}a_{i\ell}b_{\ell j}.
\]

\begin{definition}
Транспонированная матрица $A^T$ задаётся равенством
\[
  (a_{ij})^T=(a_{ji}).
\]
\end{definition}

\begin{definition}
След квадратной матрицы --- сумма элементов главной диагонали:
\[
  \tr A=a_{11}+a_{22}+\cdots+a_{nn}.
\]
\end{definition}

Имеют место равенства
\[
  \tr(A+B)=\tr A+\tr B,
  \qquad
  \tr(AB)=\tr(BA).
\]

Множество $M_n(F)$ с введёнными операциями является ассоциативным кольцом
с единицей. В общем случае оно некоммутативно, имеет делители нуля, а
уравнение $AX=E$ при $A\ne0$ не всегда имеет решение.


Если размеры блоков согласованы, сложение и умножение блочных матриц
выполняются по обычным правилам.

\begin{definition}
Матрица называется симметрической, если $A=A^T$, и кососимметрической,
если $A=-A^T$. Две матрицы называются коммутирующими, если $AB=BA$.
\end{definition}

\section{Элементарные матрицы}

Обозначим через $E_{ij}$ матричную единицу: в позиции $(i,j)$ стоит
единица, а на остальных местах --- нули. Символ Кронекера
\[
  \delta_{ij}=\begin{cases}1,&i=j,\\0,&i\ne j\end{cases}
\]
даёт формулу
\[
  E_{ij}E_{k\ell}=\delta_{jk}E_{i\ell}.
\]

Три типа элементарных матриц:
\begin{enumerate}[label=\arabic*)]
  \item растяжение строки $i$:
  \[
    T_i(\lambda)=E+(\lambda-1)E_{ii};
  \]
  \item перестановка строк $i$ и $j$:
  \[
    P_{ij}=E-E_{ii}-E_{jj}+E_{ij}+E_{ji};
  \]
  \item прибавление к строке $i$ строки $j$, умноженной на $\lambda$:
  \[
    D_{ij}(\lambda)=E+\lambda E_{ij}.
  \]
\end{enumerate}
Умножение матрицы слева на соответствующую элементарную матрицу выполняет
указанное преобразование строк.


\chapter{Евклидовы кольца}

\begin{definition}
Область целостности --- ассоциативное коммутативное кольцо с единицей и
без делителей нуля.
\end{definition}

\begin{definition}
Евклидово кольцо --- область целостности $R$, для которой существует
функция
\[
  N\colon R\setminus\{0\}\to\N_0
\]
со свойствами:
\begin{enumerate}[label=\arabic*)]
  \item $N(fg)\geq N(f)$; равенство $N(fg)=N(f)$ имеет место тогда и
  только тогда, когда $g$ обратим;
  \item для любых $f,g\in R$, $g\ne0$, существуют $q,r\in R$ такие, что
  \[
    f=gq+r,
  \]
  причём либо $r=0$, либо $N(r)<N(g)$.
\end{enumerate}
\end{definition}

В $\Z$ можно взять $N(a)=|a|$; например,
$17=5\cdot3+2$ и $N(2)<N(5)$. В кольце многочленов над полем роль нормы
играет степень.

\begin{example}
\[
  x^3+2x+1=(x^2-1)(x+1)+(2x+2),
\]
причём $\deg(2x+2)=1<2=\deg(x^2-1)$.
\end{example}

\begin{definition}
Неприводимый элемент --- элемент, который нельзя разложить в произведение
необратимых элементов.
\end{definition}

\begin{definition}
Элементы называются взаимно простыми, если их наибольший общий делитель
обратим. Простой элемент --- необратимый ненулевой элемент области
целостности, который нельзя представить произведением двух необратимых
элементов.
\end{definition}

\begin{editorialnote}
Последняя фраза дословно описывает \emph{неприводимый} элемент. Обычно
простым называют элемент $p$, для которого
$p\mid ab\Rightarrow p\mid a$ или $p\mid b$. В произвольной области
целостности эти понятия могут различаться, но в евклидовом, а значит
факториальном, кольце они совпадают.
\end{editorialnote}

\begin{definition}
Область целостности называется факториальной, если:
\begin{enumerate}[label=\arabic*)]
  \item каждый ненулевой необратимый элемент раскладывается в произведение
  неприводимых;
  \item такое разложение единственно с точностью до порядка множителей и
  умножения на обратимые элементы.
\end{enumerate}
\end{definition}


\begin{theorem}
Любое евклидово кольцо факториально.
\end{theorem}

\begin{proof}
Сначала евклидов алгоритм даёт тождество Безу: для любых $a,b\in R$
их наибольший общий делитель представим в виде $ua+vb$.

Докажем существование разложения по индукции по $N(a)$. Если $a$
неприводим, всё готово. Иначе $a=bc$, где $b$ и $c$ необратимы.
Из свойства нормы следует $N(b)<N(a)$ и $N(c)<N(a)$, поэтому оба
множителя раскладываются в произведения неприводимых.

Покажем, что каждый неприводимый $p$ прост. Если $p\mid ab$ и
$p\nmid a$, то $\gcd(p,a)$ обратим. По тождеству Безу
$up+va=1$. Умножая на $b$, получаем
$upb+vab=b$; оба слагаемых слева делятся на $p$, значит, $p\mid b$.

Теперь единственность разложения следует стандартно: неприводимый
множитель одной стороны делит произведение на другой стороне, а потому
ассоциирован с одним из её множителей. Сокращая их, продолжаем
индукцией.
\end{proof}

\chapter{Идеалы и факторкольца}

\begin{definition}
Аддитивная подгруппа $I$ кольца $R$ называется левым идеалом, если
\[
  r\in R,\ x\in I\Rightarrow rx\in I,
\]
и правым идеалом, если
\[
  r\in R,\ x\in I\Rightarrow xr\in I.
\]
Если левый и правый идеалы совпадают, говорят о двустороннем идеале.
\end{definition}

\begin{definition}
Главный идеал, порождённый элементом $a$, имеет вид
\[
  (a)=\{ra\mid r\in R\}
\]
в коммутативном кольце (для односторонних идеалов порядок множителей
учитывается).
\end{definition}

\begin{definition}
Кольцо, в котором все идеалы главные, называется кольцом главных идеалов.
\end{definition}

\begin{definition}
Пусть $I$ --- двусторонний идеал кольца $R$. На $R$ зададим отношение
\[
  a\sim b\quad\Longleftrightarrow\quad a-b\in I.
\]
Множество классов $R/I$ с операциями
\[
  (a+I)+(b+I)=(a+b)+I,\qquad
  (a+I)(b+I)=ab+I
\]
называется факторкольцом $R$ по идеалу $I$.
\end{definition}

\begin{proposition}
Операции в $R/I$ не зависят от выбора представителей классов, а
каноническая проекция
\[
  \pi\colon R\to R/I,\qquad a\mapsto a+I,
\]
является сюръективным гомоморфизмом колец с ядром $I$.
\end{proposition}

\begin{proof}
Если $a-a'\in I$ и $b-b'\in I$, то
\[
  (a+b)-(a'+b')\in I.
\]
Кроме того,
\[
  ab-a'b'=a(b-b')+(a-a')b'\in I,
\]
поскольку $I$ двусторонний. Значит, оба результата зависят только от
классов. Свойства $\pi$ следуют непосредственно из определения, а
$\pi(a)=I$ равносильно $a\in I$.
\end{proof}

\begin{explanation}
Требование двусторонности существенно именно для умножения классов.
Для коммутативного кольца всякий идеал автоматически двусторонний.
\end{explanation}


\chapter{Определитель и его приложения}

\begin{definition}
Определитель --- функция
\[
  \det\colon F^n\times\cdots\times F^n\to F
\]
от $n$ строк (или столбцов), удовлетворяющая условиям:
\begin{enumerate}[label=\arabic*)]
  \item полилинейность по каждому аргументу, например
  \[
    \det(\ldots,\alpha v+\beta w,\ldots)
    =\alpha\det(\ldots,v,\ldots)
     +\beta\det(\ldots,w,\ldots);
  \]
  \item альтернированность: если два аргумента равны, то определитель
  равен нулю;
  \item нормировка:
  \[
    \det(e_1,e_2,\ldots,e_n)=1.
  \]
\end{enumerate}
\end{definition}

Из полилинейности и альтернированности следует кососимметричность:
\[
  \det(\ldots,v,\ldots,w,\ldots)
  =-\det(\ldots,w,\ldots,v,\ldots).
\]
Действительно, раскрыв
$\det(\ldots,v+w,\ldots,v+w,\ldots)=0$, получаем, что два смешанных
слагаемых имеют противоположные знаки. В характеристике $2$ знак ``минус''
совпадает со знаком ``плюс'', но утверждение остаётся верным.


\begin{theorem}[Существование и единственность определителя]
Существует единственная функция с перечисленными свойствами, и она задаётся
формулой Лейбница
\[
  \det A=\sum_{\sigma\in S_n}
  \sgn(\sigma)\,
  a_{\sigma(1),1}a_{\sigma(2),2}\cdots a_{\sigma(n),n}.
\]
\end{theorem}

В каждом произведении выбирается по одному элементу из каждого столбца и
каждой строки. Для $n=2$
\[
  \det\begin{pmatrix}a_{11}&a_{12}\\a_{21}&a_{22}\end{pmatrix}
  =a_{11}a_{22}-a_{21}a_{12}.
\]
Полилинейность следует из распределения суммы по выбранному столбцу;
остальные свойства проверяются непосредственно по перестановкам.

\begin{proof}
Пусть столбцы матрицы равны
$v_j=\sum_{i=1}^{n}a_{ij}e_i$. По полилинейности любой определитель с
заданными свойствами обязан удовлетворять
\[
  \det(v_1,\ldots,v_n)
  =\sum_{i_1,\ldots,i_n}
   a_{i_1,1}\cdots a_{i_n,n}
   \det(e_{i_1},\ldots,e_{i_n}).
\]
Если два индекса совпадают, соответствующее слагаемое равно нулю.
Поэтому остаются только наборы
$(i_1,\ldots,i_n)=(\sigma(1),\ldots,\sigma(n))$. Переставляя базисные
векторы в стандартный порядок, получаем
\[
  \det(e_{\sigma(1)},\ldots,e_{\sigma(n)})=\sgn(\sigma).
\]
Это доказывает единственность и приводит к формуле Лейбница.
Обратно, правая часть формулы непосредственно полилинейна,
альтернирована и на $(e_1,\ldots,e_n)$ равна $1$, значит, задаёт
требуемую функцию.
\end{proof}


Из аксиом и формулы Лейбница следуют правила:
\begin{enumerate}[label=\arabic*)]
  \item $\det(A T_i(\lambda))=\lambda\det A$;
  \item $\det(A D_{ij}(\lambda))=\det A$;
  \item $\det(A P_{ij})=-\det A$;
  \item $\det A=\det A^T$;
  \item определитель треугольной матрицы равен произведению диагональных
  элементов;
  \item для блочно-треугольной матрицы
  \[
    \det\begin{pmatrix}A&*\\0&B\end{pmatrix}
    =\det A\det B.
  \]
\end{enumerate}

\begin{definition}
Минор $M_{ij}$ --- определитель матрицы, полученной вычёркиванием
$i$-й строки и $j$-го столбца. Алгебраическое дополнение
\[
  A_{ij}=(-1)^{i+j}M_{ij}.
\]
\end{definition}

Разложение определителя по $i$-й строке или $j$-му столбцу имеет вид
\[
  \det A=\sum_{j=1}^{n}a_{ij}A_{ij}
  =\sum_{i=1}^{n}a_{ij}A_{ij}.
\]


\section{Системы линейных уравнений}

Квадратную систему
\[
\begin{cases}
a_{11}x_1+a_{12}x_2+\cdots+a_{1n}x_n=b_1,\\
a_{21}x_1+a_{22}x_2+\cdots+a_{2n}x_n=b_2,\\
\hfill\vdots\\
a_{n1}x_1+a_{n2}x_2+\cdots+a_{nn}x_n=b_n
\end{cases}
\]
записывают как $Ax=b$, где $A\in M_n(F)$, $x$ --- столбец неизвестных,
$b$ --- столбец свободных членов.

\begin{theorem}[Крамера]
Если $\det A\ne0$, то система $Ax=b$ имеет единственное решение
\[
  x_i=\frac{\det A_i}{\det A},
\]
где $A_i$ получена из $A$ заменой $i$-го столбца столбцом $b$.
\end{theorem}

\begin{proof}
Пусть $a_1,\ldots,a_n$ --- столбцы $A$. Из $Ax=b$ следует
$b=x_1a_1+\cdots+x_na_n$. Заменим в определителе $i$-й столбец на $b$:
\[
  \det A_i
  =\det(a_1,\ldots,b,\ldots,a_n)
  =x_i\det A.
\]
Все остальные слагаемые исчезают, поскольку содержат два одинаковых
столбца. Отсюда получается формула Крамера. Единственность следует из
обратимости $A$ при $\det A\ne0$.
\end{proof}

\begin{example}
Для системы
\[
  \begin{cases}x+y=2,\\x-y=0\end{cases}
\]
имеем
\[
  \det\begin{pmatrix}1&1\\1&-1\end{pmatrix}=-2\ne0,
  \qquad x=y=1.
\]
\end{example}

\begin{definition}
Матрица $A^{-1}$ называется обратной к $A$, если
\[
  AA^{-1}=A^{-1}A=E.
\]
\end{definition}

Матрица обратима тогда и только тогда, когда $\det A\ne0$; при этом
\[
  \det A^{-1}=\frac1{\det A},
  \qquad
  A^{-1}=\frac1{\det A}
  \begin{pmatrix}
  A_{11}&A_{12}&\cdots&A_{1n}\\
  A_{21}&A_{22}&\cdots&A_{2n}\\
  \vdots&\vdots&\ddots&\vdots\\
  A_{n1}&A_{n2}&\cdots&A_{nn}
  \end{pmatrix}^{T}.
\]

\begin{explanation}
Матрица в числителе называется присоединённой и обозначается
$\operatorname{adj}A$. Разложение определителя по строке и столбцу даёт
\[
  A\,\operatorname{adj}A
  =\operatorname{adj}A\,A
  =(\det A)E.
\]
При $\det A\ne0$ деление на $\det A$ даёт формулу для $A^{-1}$.
Обратно, если $A$ обратима, то
$1=\det E=\det A\det A^{-1}$, поэтому $\det A\ne0$.
\end{explanation}


Дополнительные свойства:
\begin{enumerate}[label=\arabic*)]
  \item $\det(AT_i(\lambda))=\lambda\det A$;
  \item $\det(AD_{ij}(\lambda))=\det A$;
  \item $\det(AP_{ij})=-\det A$;
  \item $\det A=\det A^T$;
  \item определитель треугольной матрицы равен произведению её
  диагональных элементов;
  \item $\det\begin{pmatrix}A&*\\0&B\end{pmatrix}=\det A\det B$.
\end{enumerate}

\chapter{Модули над кольцом}

Пусть $R$ --- ассоциативное кольцо, а $M$ --- абелева группа.

\begin{definition}
Правый $R$-модуль --- множество $M$ с отображением
\[
  M\times R\to M,\qquad (x,\lambda)\mapsto x\lambda,
\]
для которого
\[
\begin{aligned}
  (x\lambda)\mu&=x(\lambda\mu),\\
  x(\lambda+\mu)&=x\lambda+x\mu,\\
  (x+y)\lambda&=x\lambda+y\lambda.
\end{aligned}
\]
Левый модуль определяется аналогично отображением $R\times M\to M$.
\end{definition}

\begin{definition}
Модуль над кольцом с единицей называется унитальным, если
$x\cdot1=x$ (соответственно $1\cdot x=x$) для всех $x\in M$.
\end{definition}

Любое кольцо является модулем над самим собой. Если $I$ --- правый
(левый) идеал кольца $R$, то $I$ --- правый (левый) $R$-модуль.

\begin{definition}
Противоположное кольцо $R^{\circ}$ имеет то же аддитивное строение, а
умножение задаётся равенством $a\circ b=ba$. Правый $R^{\circ}$-модуль
эквивалентен левому $R$-модулю.
\end{definition}

\begin{definition}
Подмодуль $N\leq M$ --- аддитивная подгруппа, замкнутая относительно
умножения на скаляры: $x\in N$, $a\in R\Rightarrow ax\in N$.
\end{definition}

\begin{editorialnote}
Начиная с этой формулы используется соглашение о \emph{левых} модулях:
скаляр записывается слева, $ax$. Для правого модуля условие следует
писать как $xa\in N$; все последующие утверждения имеют симметричный
правый вариант.
\end{editorialnote}

Фактормодуль $M/N$ состоит из смежных классов, причём
\[
  \bar m+\bar n=\overline{m+n},
  \qquad a\bar m=\overline{am}.
\]

Гомоморфизм модулей $\varphi\colon M\to N$ сохраняет сложение и
умножение на скаляры:
\[
  \varphi(x+y)=\varphi(x)+\varphi(y),
  \qquad
  \varphi(ax)=a\varphi(x).
\]


Дальнейшие определения формулируются для унитальных модулей над кольцом с
единицей.

\begin{definition}
Элементы $x_1,\ldots,x_n$ порождают модуль $M$, если каждый $v\in M$
представляется конечной линейной комбинацией
\[
  v=\lambda_1x_1+\lambda_2x_2+\cdots+\lambda_nx_n,
  \qquad \lambda_i\in R.
\]
Элементы $x_i$ называются порождающими (образующими) модуля $M$.
\end{definition}

\begin{definition}
Система порождающих называется базисом, если указанное представление
единственно для каждого $v\in M$.
\end{definition}

\begin{definition}
Свободный модуль --- модуль, обладающий базисом. Векторное пространство
всегда является свободным модулем над полем.
\end{definition}

Рассмотрим два разложения одного элемента:
\[
  v=\lambda_1x_1+\cdots+\lambda_nx_n,
  \qquad
  v=\mu_1x_1+\cdots+\mu_nx_n.
\]
Тогда
\[
  0=(\lambda_1-\mu_1)x_1+\cdots+(\lambda_n-\mu_n)x_n.
\]

\begin{definition}
Нетривиальная линейная комбинация --- комбинация, в которой не все
коэффициенты равны нулю. Элементы $x_1,\ldots,x_n$ называются линейно
зависимыми, если существует их нетривиальная линейная комбинация, равная
нулю; в противном случае они линейно независимы.
\end{definition}

Базис модуля --- линейно независимая система порождающих.


Термин ``система векторов'' используется как синоним ``множество
векторов'', при этом векторы могут занумеровываться и могут повторяться.

\begin{definition}[Критерий зависимости]
Система векторов линейно зависима тогда и только тогда, когда один из её
векторов линейно выражается через остальные.
\end{definition}

\begin{editorialnote}
Для произвольного модуля над кольцом это утверждение требует уточнения:
в разрешаемом коэффициенте должен встречаться обратимый элемент. Например,
в $\Z/4\Z$ как $\Z$-модуле ненулевой элемент $\bar2$ зависим, поскольку
$2\bar2=0$, но выразить $\bar2$ через пустую систему нельзя. Без
оговорок критерий верен для векторных пространств над полем и ниже
доказывается именно в этой форме.
\end{editorialnote}

\chapter{Векторные пространства}

\begin{definition}
Если $F$ --- поле, то $F$-векторное пространство --- унитальный
$F$-модуль: множество $V$ с операциями сложения
$V\times V\to V$ и умножения на скаляры
$V\times F\to V$, удовлетворяющими аксиомам модуля.
\end{definition}

\begin{definition}
Алгеброй над полем $F$ называется векторное пространство $A$ с
дополнительной операцией умножения $A\times A\to A$, дистрибутивной по
сложению и совместимой с умножением на скаляры:
\[
  \lambda(ab)=(\lambda a)b=a(\lambda b).
\]
\end{definition}

Алгебра объединяет свойства векторного пространства и кольца. Подпространство
$U\leq V$ замкнуто относительно сложения и умножения на скаляры.

\begin{definition}
Подалгебра $B$ алгебры $A$ является одновременно подкольцом алгебры $A$
и подпространством её векторного пространства.
\end{definition}


\begin{theorem}[Критерий линейной зависимости]
Система векторов $v_1,\ldots,v_n$ линейно зависима тогда и только тогда,
когда хотя бы один вектор линейно выражается через остальные.
\end{theorem}

\begin{proof}
Если
\[
  v_1\lambda_1+\cdots+v_n\lambda_n=0
\]
и некоторый $\lambda_i\ne0$, то в векторном пространстве можно умножить
на $\lambda_i^{-1}$ и получить
\[
  v_i=-\sum_{j\ne i}v_j\lambda_j\lambda_i^{-1}.
\]
Обратно, перенеся выражение одного вектора через остальные в одну сторону,
получаем нетривиальную линейную комбинацию, равную нулю.
\end{proof}

\begin{definition}
Линейная оболочка векторов $a_1,\ldots,a_n\in V$ --- множество всех их
линейных комбинаций:
\[
  \langle a_1,\ldots,a_n\rangle
  =L(a_1,\ldots,a_n)
  =\left\{\sum_{i=1}^{n}a_i\lambda_i\,\middle|\,\lambda_i\in F\right\}
  \leq V.
\]
Если эта оболочка равна $V$, то векторы порождают $V$.
\end{definition}


\chapter{Системы линейных алгебраических уравнений}

Пусть $a_1,\ldots,a_n,b\in F^k$. Векторы $a_i$ порождают некоторое
подпространство. Вопрос о принадлежности $b$ этому подпространству
эквивалентен представлению
\[
  b=a_1x_1+\cdots+a_nx_n.
\]
Покоординатно это система
\[
\begin{cases}
a_{11}x_1+a_{12}x_2+\cdots+a_{1n}x_n=b_1,\\
a_{21}x_1+a_{22}x_2+\cdots+a_{2n}x_n=b_2,\\
\hfill\vdots\\
a_{k1}x_1+a_{k2}x_2+\cdots+a_{kn}x_n=b_k.
\end{cases}
\]
В матричном виде: $Ax=b$, где $A\in M_{k\times n}(F)$, а решение ---
столбец $x\in F^n$.

\begin{definition}
Система называется совместной, если существует хотя бы одно решение;
определённой, если решение единственно; несовместной, если решений нет.
Эквивалентные системы имеют одинаковое множество решений.
\end{definition}


\begin{definition}
Ведущий элемент строки --- её первый ненулевой элемент. Ступенчатая
матрица --- матрица, у которой положения ведущих элементов строк образуют
строго возрастающую последовательность, а все нулевые строки находятся в
конце.
\end{definition}

\begin{theorem}
Любую матрицу можно привести к ступенчатому виду элементарными
преобразованиями строк.
\end{theorem}

\begin{proof}
Алгоритм:
\begin{enumerate}[label=\Roman*.]
  \item Если матрица нулевая, она уже ступенчатая.
  \item Иначе найти первый ненулевой столбец.
  \item Переставить строки так, чтобы ненулевой элемент этого столбца
  оказался в первой строке, затем вычесть подходящие кратные первой строки
  из последующих, зануляя остальные элементы столбца.
  \item Повторить процедуру для оставшейся подматрицы.
\end{enumerate}
После каждого шага либо уменьшается число рассматриваемых строк, либо
число рассматриваемых столбцов. Поэтому процесс конечен и заканчивается
ступенчатой матрицей.
\end{proof}

Это метод Гаусса. Он заключается в приведении расширенной матрицы системы
$[A\mid b]$ к ступенчатому виду. Пусть $r$ --- число ненулевых строк в
$A$, $\widetilde r$ --- число ненулевых строк в $[A\mid b]$, а $n$ ---
число неизвестных.


Возможны три случая.
\begin{enumerate}[label=\arabic*)]
  \item $r<\widetilde r$: система несовместна, поскольку появляется
  строка $0x_1+\cdots+0x_n=c$ с $c\ne0$.
  \item $r=\widetilde r=n$: система определённа. Обратной подстановкой
  последовательно находятся $x_n,x_{n-1},\ldots,x_1$.
  \item $r=\widetilde r<n$: система совместна и имеет свободные
  переменные. Ведущие переменные выражаются через свободные, а свободным
  придаются произвольные значения из поля.
\end{enumerate}

Вывод: свободные переменные параметризуют все решения; если поле имеет
более одного элемента, наличие свободной переменной даёт более одного
решения. Метод Гаусса не только определяет совместность и число решений,
но и позволяет найти их.

Прямой ход Гаусса --- приведение к ступенчатому виду. Обратный ход ---
последовательное нахождение ведущих неизвестных во втором случае.


\begin{example}
Система
\[
  \begin{cases}x-y=0,\\x+y=2\end{cases}
\]
имеет единственное решение $(1,1)$, поскольку
\[
  \left(\begin{array}{cc|c}1&-1&0\\1&1&2\end{array}\right)
  \sim
  \left(\begin{array}{cc|c}1&-1&0\\0&2&2\end{array}\right).
\]
\end{example}

\begin{example}
Система
\[
  \begin{cases}x-y=1,\\2x-2y=2\end{cases}
\]
имеет бесконечно много решений: $x=1+y$, $y\in F$. Геометрически обе
прямые совпадают.
\end{example}

\begin{example}
Система
\[
  \begin{cases}x-y=1,\\2x-2y=4\end{cases}
\]
несовместна: после исключения получается $0=2$. Геометрически прямые
параллельны.
\end{example}

\begin{definition}
Система $Ax=0$ называется однородной, а $Ax=b$ при $b\ne0$ ---
неоднородной.
\end{definition}

Однородная система всегда совместна и всегда имеет нулевое решение.
Если однородная система с числом уравнений, меньшим числа неизвестных,
приведена к ступенчатому виду, то она имеет ненулевое решение.


Совокупность всех решений однородной системы $Ax=0$ с $n$ неизвестными
является подпространством в $F^n$: если $Ax=Ay=0$, то
\[
  A(\lambda x+\mu y)=\lambda Ax+\mu Ay=0.
\]

\begin{theorem}[О геометрической структуре решений]
Если неоднородная система $Ax=b$ совместна и $x_{\mathrm{ч}}$ --- одно
частное решение, а $X_0=\{x\mid Ax=0\}$ --- пространство решений
однородной системы, то множество всех решений равно
\[
  X=X_0+x_{\mathrm{ч}}.
\]
\end{theorem}

\begin{proof}
Для $x_0\in X_0$ имеем $A(x_{\mathrm{ч}}+x_0)=b$. Если $x$ и
$x_{\mathrm{ч}}$ --- два решения неоднородной системы, то
$A(x-x_{\mathrm{ч}})=0$, следовательно, $x-x_{\mathrm{ч}}\in X_0$.
\end{proof}

\begin{definition}
Линейное многообразие --- подмножество пространства, полученное сдвигом
подпространства на фиксированный вектор.
\end{definition}


\chapter{Базис и размерность векторных пространств}

\begin{theorem}[Основная лемма линейной зависимости]
Если векторы $b_1,\ldots,b_m$ выражаются через
$a_1,\ldots,a_n$ и $m>n$, то система $b_1,\ldots,b_m$ линейно зависима.
\end{theorem}

\begin{proof}
Запишем
\[
  b_j=\sum_{i=1}^{n}a_i c_{ij},\qquad j=1,\ldots,m,
\]
и составим матрицу коэффициентов $C=(c_{ij})\in M_{n\times m}(F)$.
Однородная система $C\lambda=0$ имеет $m>n$ неизвестных, поэтому после
приведения к ступенчатому виду остаётся хотя бы одна свободная переменная.
Следовательно, существует ненулевой столбец $\lambda$ такой, что
\[
  \sum_{j=1}^{m}\lambda_jb_j
  =\sum_{i=1}^{n}a_i\left(\sum_{j=1}^{m}c_{ij}\lambda_j\right)=0.
\]
Это нетривиальная линейная зависимость между $b_1,\ldots,b_m$.
\end{proof}

\begin{definition}
Конечномерное пространство --- пространство, порождённое конечным числом
векторов.
\end{definition}

\begin{definition}
Базис конечномерного пространства $V$ --- линейно независимая система
векторов, порождающая $V$.
\end{definition}

\begin{theorem}[Существование базиса]
В любом конечномерном пространстве существует базис.
\end{theorem}

\begin{proof}
Если система порождающих линейно зависима, один из векторов выражается
через остальные и может быть удалён без изменения линейной оболочки.
Повторяя процесс, получаем базис.
\end{proof}

\begin{theorem}[О числе векторов в базисе]
Все базисы конечномерного векторного пространства содержат одинаковое
число векторов.
\end{theorem}


\begin{proof}
Пусть $e_1,\ldots,e_n$ и $\widetilde e_1,\ldots,\widetilde e_m$ --- два
базиса. Каждый вектор первого базиса выражается через второй, поэтому
$n\leq m$. Меняя базисы местами, получаем $m\leq n$, следовательно,
$m=n$.
\end{proof}

\begin{definition}
Размерность конечномерного пространства $V$ --- число векторов в любом
его базисе; обозначение $\dim V$.
\end{definition}

\begin{theorem}[Дополнение до базиса]
Любую линейно независимую систему в конечномерном пространстве можно
дополнить до базиса этого пространства.
\end{theorem}

\begin{proof}
Пусть $v_1,\ldots,v_k$ линейно независимы. Если они уже порождают $V$,
то это базис. Иначе выберем
$v_{k+1}\notin\langle v_1,\ldots,v_k\rangle$; новая система остаётся
линейно независимой. Продолжаем этот процесс. Он завершится не позднее
чем после $\dim V-k$ шагов, поскольку линейно независимая система не
может содержать больше $\dim V$ векторов.
\end{proof}

\begin{theorem}[Монотонность размерности]
Если $U\subseteq V$ --- подпространство конечномерного пространства, то
\[
  \dim U\leq\dim V.
\]
\end{theorem}

\begin{proof}
Базис $U$ является линейно независимой системой в $V$ и потому
дополняется до базиса $V$. Следовательно, число его элементов не
превосходит $\dim V$.
\end{proof}

\begin{theorem}
Если $V$ и $U$ --- конечномерные пространства над $F$, причём
$\dim V=\dim U=n$, то
\[
  V\cong F^n,
  \qquad
  V\cong U.
\]
\end{theorem}

\begin{proof}
Выберем базис $e_1,\ldots,e_n$ пространства $V$. Координатное
отображение
\[
  x_1e_1+\cdots+x_ne_n\longmapsto(x_1,\ldots,x_n)
\]
является изоморфизмом $V\cong F^n$. Аналогично $U\cong F^n$, поэтому
$V\cong U$.
\end{proof}


\chapter{Ранг матрицы}

Для матрицы $A$ равносильны следующие определения ранга:
\begin{enumerate}[label=\arabic*)]
  \item ранг системы её строк, то есть размерность линейной оболочки
  строк;
  \item строчный ранг --- размерность линейной оболочки строк матрицы
  $A\in M_{k\times n}(F)$;
  \item столбцовый ранг --- размерность линейной оболочки столбцов;
  \item минорный ранг --- наибольший порядок ненулевого минора матрицы.
\end{enumerate}

Строки и столбцы, соответствующие базисному минору, линейно независимы и
называются базисными. Остальные строки (столбцы) выражаются через базисные
линейными комбинациями.


\begin{theorem}
Строчный, столбцовый и минорный ранги матрицы совпадают:
\[
  \operatorname{rowrank}A
  =\operatorname{colrank}A
  =\rk A.
\]
\end{theorem}

\begin{proof}
Элементарными преобразованиями строк и столбцов матрицу можно привести
к нормальной форме
\[
  PAQ=
  \begin{pmatrix}
    E_r&0\\0&0
  \end{pmatrix},
\]
где $P$ и $Q$ обратимы. Обратимые преобразования не изменяют размерности
оболочек строк и столбцов и не меняют наибольший порядок ненулевого
минора. В нормальной форме все три числа равны $r$, следовательно, они
равны и для исходной матрицы.
\end{proof}

Для $A\in M_n(F)$ условие $\rk A=n$ равносильно $\det A\ne0$.

\begin{theorem}
Ранг матрицы не изменяется при элементарных преобразованиях строк.
\end{theorem}

\begin{proof}
Каждое элементарное преобразование строк равносильно умножению слева на
обратимую элементарную матрицу $E$. Поэтому
\[
  \Im(EA)=E(\Im A),
\]
а обратимое линейное отображение $E$ сохраняет размерность образа.
\end{proof}

\begin{theorem}[Кронекера--Капелли]
Система $Ax=b$ совместна тогда и только тогда, когда
\[
  \rk A=\rk[A\mid b].
\]
\end{theorem}

\begin{proof}
Система совместна тогда и только тогда, когда $b$ является линейной
комбинацией столбцов $A$. Это равносильно тому, что добавление столбца
$b$ не увеличивает размерность их линейной оболочки, то есть
$\rk A=\rk[A\mid b]$.
\end{proof}

\begin{theorem}
Пространство решений однородной системы $Ax=0$ с $n$ неизвестными имеет
размерность
\[
  n-\rk A.
\]
\end{theorem}

\begin{proof}
После приведения $A$ к ступенчатому виду число ведущих переменных равно
$\rk A$, а число свободных --- $n-\rk A$. Каждому набору значений
свободных переменных соответствует единственное решение, поэтому
пространство решений изоморфно $F^{\,n-\rk A}$.
\end{proof}


\chapter{Гомоморфизмы векторных пространств}

Пусть $U,V$ --- векторные пространства над полем $F$.

\begin{definition}
Гомоморфизм (линейное отображение) $\varphi\colon U\to V$ удовлетворяет
\[
  \varphi(\alpha a+\beta b)
  =\alpha\varphi(a)+\beta\varphi(b)
\]
для всех $a,b\in U$ и $\alpha,\beta\in F$. В частности,
\[
  \varphi(0)=0,
  \quad \varphi(-a)=-\varphi(a),
  \quad \varphi(a-b)=\varphi(a)-\varphi(b).
\]
\end{definition}

Множество всех линейных отображений обозначается
$\operatorname{Hom}(U,V)$, а
$\operatorname{Hom}(U,U)=\operatorname{End}(U)$.


Пусть $\dim U=m$, $\dim V=n$, $(e_1,\ldots,e_m)$ --- базис $U$, а
$(f_1,\ldots,f_n)$ --- базис $V$. Для любого $x\in U$ его координаты в
первом базисе определяют разложение, а значения отображения на базисных
векторах имеют вид
\[
  \varphi(e_j)=\sum_{i=1}^{n}a_{ij}f_i.
\]
Коэффициенты $a_{ij}$ составляют матрицу линейного отображения
$A\in M_{n\times m}(F)$; её столбцы --- координаты образов базисных
векторов. Если $X$ и $Y$ --- столбцы координат $x$ и $\varphi(x)$, то
\[
  Y=AX.
\]
Таким образом, уравнение $Ax=b$ спрашивает, является ли $b$ образом
некоторого вектора при гомоморфизме с матрицей $A$.

\begin{definition}
Образ и ядро линейного отображения $A\colon U\to V$:
\[
  \Im A=\{y\in V\mid \exists x\in U:\ A(x)=y\}\leq V,
\]
\[
  \Ker A=\{x\in U\mid A(x)=0\}\leq U.
\]
Ранг отображения --- $\dim\Im A$.
\end{definition}

\begin{theorem}[О ранге и дефекте]
Для $A\colon U\to V$ при $\dim U=n$
\[
  \dim\Im A+\dim\Ker A=\dim U=n.
\]
\end{theorem}

\begin{proof}
Выберем базис $e_1,\ldots,e_k$ ядра $\Ker A$ и дополним его до базиса
$U$:
\[
  e_1,\ldots,e_k,e_{k+1},\ldots,e_n.
\]
Тогда векторы $Ae_{k+1},\ldots,Ae_n$ образуют базис $\Im A$.
Они порождают образ, поскольку первые $k$ базисных векторов переходят в
нуль; их линейная независимость следует из того, что линейная комбинация
$e_{k+1},\ldots,e_n$, попавшая в ядро, должна лежать одновременно в
оболочке этих векторов и в оболочке $e_1,\ldots,e_k$. Следовательно,
\[
  \dim\Ker A=k,\qquad \dim\Im A=n-k.
\]
\end{proof}


Для матрицы отображения
\[
  \dim\Im A=\rk A,
\]
поскольку образ порождается столбцами матрицы. Поэтому предыдущая теорема
даёт
\[
  \dim\Im A+\dim\Ker A=\dim U.
\]

На $\operatorname{Hom}(U,V)$ вводятся операции
\[
  (A+B)(x)=A(x)+B(x),
  \qquad
  (\lambda A)(x)=\lambda A(x).
\]
Они превращают $\operatorname{Hom}(U,V)$ в векторное пространство
размерности
\[
  \dim U\cdot\dim V,
\]
а выбор базисов отождествляет линейные отображения с матрицами:
$A,B\leftrightarrow A+B$ и $\lambda A\leftrightarrow\lambda A$.

Композиции отображений соответствует произведение матриц. Если
$V\xrightarrow{B}U\xrightarrow{A}W$, то
\[
  [A\circ B]=[A][B].
\]

\begin{proposition}
Если $\dim U=\dim V<\infty$, то для линейного отображения
$\varphi\colon U\to V$ равносильны условия:
\begin{enumerate}[label=\arabic*)]
  \item $\varphi$ --- изоморфизм;
  \item $\Ker\varphi=\{0\}$;
  \item $\Im\varphi=V$,
\end{enumerate}
\end{proposition}

\begin{proof}
Условие $\Ker\varphi=\{0\}$ по теореме о ранге и дефекте даёт
$\dim\Im\varphi=\dim U=\dim V$, поэтому $\Im\varphi=V$. Обратно,
сюръективность даёт $\dim\Im\varphi=\dim V=\dim U$, откуда
$\dim\Ker\varphi=0$. Линейное отображение является изоморфизмом тогда и
только тогда, когда оно одновременно инъективно и сюръективно.
\end{proof}

\begin{explanation}
При разных размерностях инъективность и сюръективность уже не
равносильны и рассматриваются отдельно.
\end{explanation}

\begin{visualreconstruction}
Коэффициенты матрицы композиции удобно получить одной цепочкой. При
базисах $e_i$, $f_j$, $g_\ell$:
\[
  B(e_i)=\sum_{j=1}^{m}b_{ji}f_j,
  \qquad
  A(f_j)=\sum_{\ell=1}^{n}a_{\ell j}g_\ell,
\]
\[
  (A\circ B)(e_i)
  =\sum_{j=1}^{m}\sum_{\ell=1}^{n}a_{\ell j}b_{ji}g_\ell
  =\sum_{\ell=1}^{n}
    \underbrace{\left(\sum_{j=1}^{m}a_{\ell j}b_{ji}\right)}_{c_{\ell i}}
    g_\ell .
\]
Следовательно, $c_{\ell i}=\sum_j a_{\ell j}b_{ji}$, то есть
$C=AB$.
\end{visualreconstruction}


\chapter[Факторпространства и операции над подпространствами]
{Факторпространства и операции над\\ подпространствами}

\begin{definition}
Пусть $U\leq V$. Векторы $v,\tilde v\in V$ называются эквивалентными
по модулю $U$, если
\[
  v-\tilde v\in U.
\]
Класс $v+U=\{\tilde v\in V\mid \tilde v-v\in U\}$; множество классов
$V/U$ с операциями
\[
  (v_1+U)+(v_2+U)=(v_1+v_2)+U,
  \qquad
  \lambda(v+U)=\lambda v+U
\]
называется факторпространством.
\end{definition}

Если $V$ конечномерно, то
\[
  \dim(V/U)=\dim V-\dim U
\]
--- коразмерность подпространства $U$.

\begin{definition}
Согласованный базис $V$ с подпространством $U$ --- базис, в котором
некоторая часть векторов образует базис $U$.
\end{definition}

Нулём в факторпространстве $V/U$ является само подпространство $U$.

\begin{theorem}[Первая теорема об изоморфизме]
Для линейного отображения $A\colon U\to V$
\[
  U/\Ker A\cong\Im A.
\]
\end{theorem}

\begin{proof}
Определим
\[
  \overline A\colon U/\Ker A\to\Im A,\qquad
  x+\Ker A\longmapsto Ax.
\]
Если $x-y\in\Ker A$, то $Ax=Ay$, поэтому отображение определено
корректно. Оно линейно и сюръективно по определению образа. Если
$\overline A(x+\Ker A)=0$, то $x\in\Ker A$, то есть
$x+\Ker A=\Ker A$; следовательно, ядро $\overline A$ нулевое и
$\overline A$ --- изоморфизм.
\end{proof}

\begin{visualreconstruction}
Пространство $U$ разбито на классы $x+\Ker A$; пространство $V$ справа
содержит вертикально отмеченный образ $\Im A$. Все точки одного класса
переходят в одну точку образа:
\begin{center}
\begin{tikzpicture}[>=Latex,scale=.92]
  \draw (0,0) rectangle (4,3.1);
  \node[above] at (2,3.1) {$U$};
  \draw (0.35,.55) rectangle (3.65,1.05);
  \draw (0.35,1.32) rectangle (3.65,1.82);
  \draw (0.35,2.09) rectangle (3.65,2.59);
  \node at (1.25,.80) {$x_3+\Ker A$};
  \node at (1.25,1.57) {$x_2+\Ker A$};
  \node at (1.25,2.34) {$x_1+\Ker A$};
  \draw (6,0) rectangle (8.35,3.1);
  \node[above] at (7.18,3.1) {$V$};
  \draw[very thick] (7.18,.35) -- (7.18,2.75);
  \node[right] at (7.18,2.72) {$\Im A$};
  \fill (7.18,.80) circle (1.6pt);
  \fill (7.18,1.57) circle (1.6pt);
  \fill (7.18,2.34) circle (1.6pt);
  \draw[->] (3.65,.80) -- (7.1,.80);
  \draw[->] (3.65,1.57) -- (7.1,1.57);
  \draw[->] (3.65,2.34) -- (7.1,2.34);
  \node[above] at (5.25,2.72) {$A$};
\end{tikzpicture}
\end{center}
Индуцированное отображение классов в $\Im A$ биективно и линейно.
\end{visualreconstruction}


Из предыдущей теоремы вновь получаем
\[
  \dim U-\dim\Ker A=\dim\Im A.
\]

Для подпространств $U,V\leq W$ определяются:
\begin{enumerate}[label=\arabic*)]
  \item пересечение $U\cap V$;
  \item объединение множеств $U\cup V$; оно обычно не является
  подпространством;
  \item сумма подпространств
  \[
    U+V=\{x+y\mid x\in U,\ y\in V\}=\langle U\cup V\rangle\leq W.
  \]
\end{enumerate}

\begin{theorem}
\[
  (U+V)/U\cong V/(U\cap V).
\]
\end{theorem}

\begin{proof}
Рассмотрим линейное отображение
\[
  \varphi\colon V\to(U+V)/U,\qquad v\mapsto v+U.
\]
Оно сюръективно: класс $(u+v)+U$ совпадает с $v+U$. Его ядро состоит из
векторов $v\in V$, лежащих в $U$, то есть равно $U\cap V$. По первой
теореме об изоморфизме
$V/(U\cap V)\cong(U+V)/U$.
\end{proof}

\begin{theorem}[Формула Грассмана]
\[
  \dim(U+V)+\dim(U\cap V)=\dim U+\dim V.
\]
\end{theorem}

\begin{proof}
Из предыдущего изоморфизма и формулы размерности факторпространства
получаем
\[
  \dim(U+V)-\dim U=\dim V-\dim(U\cap V).
\]
Перенос слагаемых даёт требуемое равенство.
\end{proof}

Следствия:
\[
  \dim(U+V)\leq\dim U+\dim V,
\]
\[
  \rk(A+B)\leq\rk A+\rk B,
\]
\[
  \rk(AB)\leq\min\{\rk A,\rk B\},
\]
поскольку
\[
  \Im(A+B)\subseteq\Im A+\Im B,\qquad
  \Im(AB)\subseteq\Im A.
\]
Кроме того, $AB$ является образом ограничения $A|_{\Im B}$, поэтому
$\rk(AB)\leq\rk B$. Вместе эти включения дают оценку через минимум.


\begin{definition}
Подпространства $U_1,\ldots,U_k$ называются линейно независимыми, если из
\[
  v_1+\cdots+v_k=0,
  \qquad v_i\in U_i,
\]
следует $v_1=\cdots=v_k=0$.
\end{definition}

При выборе базиса каждого $U_i$ их объединение линейно независимо тогда
и только тогда, когда подпространства $U_i$ линейно независимы. В этом
случае
\[
  \dim(U_1+\cdots+U_k)=\dim U_1+\cdots+\dim U_k.
\]

Для двух подпространств $U,W\leq V$ линейная независимость равносильна
$U\cap W=\{0\}$.

\begin{definition}
Пространство $V$ раскладывается во внутреннюю прямую сумму подпространств
$U_1,\ldots,U_k$, если
\begin{enumerate}[label=\arabic*)]
  \item $V=U_1+\cdots+U_k$;
  \item подпространства $U_1,\ldots,U_k$ линейно независимы.
\end{enumerate}
Обозначение:
\[
  V=U_1\oplus\cdots\oplus U_k.
\]
\end{definition}


\begin{definition}
Внешняя прямая сумма пространств $U_1,\ldots,U_k$ --- пространство
\[
  \bigoplus_{i=1}^{k}U_i
  =\{(v_1,\ldots,v_k)\mid v_i\in U_i\}
\]
с покомпонентными операциями.
\end{definition}

Внутренняя прямая сумма изоморфна внешней: элемент
$v_1+\cdots+v_k$ однозначно соответствует кортежу
$(v_1,\ldots,v_k)$.

\begin{proof}
Отображение
\[
  \Phi\colon\bigoplus_{i=1}^{k}U_i\to V,\qquad
  (v_1,\ldots,v_k)\mapsto v_1+\cdots+v_k
\]
линейно. Оно сюръективно, поскольку $V=U_1+\cdots+U_k$, и инъективно
в силу линейной независимости подпространств: равенство суммы нулю
влечёт $v_1=\cdots=v_k=0$. Следовательно, $\Phi$ --- изоморфизм.
\end{proof}

\chapter{Линейные функционалы}

Множество $\operatorname{Hom}(U,F)$ линейных отображений из $U$ в поле
$F$ называется сопряжённым (двойственным) пространством к $U$ и
обозначается $U^*$. Так как $F$ одномерно,
\[
  \dim U^*=\dim U.
\]
Линейный функционал $\omega\in U^*$ при выборе базиса
$e_1,\ldots,e_n$ записывается строкой:
\[
  x=e_1x_1+\cdots+e_nx_n,
  \qquad
  \omega(x)=\omega(e_1)x_1+\cdots+\omega(e_n)x_n.
\]
Два функционала равны тогда и только тогда, когда они совпадают на всех
векторах (достаточно проверить их значения на базисе).


Если $\omega=a_1x_1+\cdots+a_nx_n$, то коэффициенты $a_i$ задают
функционал в выбранном базисе; это однородный многочлен первой степени.
Линейная строка --- способ записи линейного функционала в фиксированном
базисе над полем.

\begin{definition}
Для базиса $e_1,\ldots,e_n$ координатные функционалы
$e_1^*,\ldots,e_n^*\in U^*$ определяются равенствами
\[
  e_i^*(x)=x_i
\]
для $x=e_1x_1+\cdots+e_nx_n$.
\end{definition}

\begin{theorem}
Координатные функционалы $e_1^*,\ldots,e_n^*$ образуют базис $U^*$.
\end{theorem}

\begin{proof}
Если $\lambda_1e_1^*+\cdots+\lambda_ne_n^*=0$, то применение к $e_i$
даёт $\lambda_i=0$, поскольку
\[
  e_j^*(e_i)=\delta_{ij}.
\]
Размерности совпадают, следовательно, линейно независимая система из $n$
элементов является базисом.
\end{proof}

\begin{definition}
$e_1^*,\ldots,e_n^*$ --- сопряжённый (двойственный) базис к
$e_1,\ldots,e_n$. Равенство $e_i^*(e_j)=\delta_{ij}$ использует символ
Кронекера.
\end{definition}

Симметрические матрицы образуют подпространство
$\{A\in M_n(F)\mid A=A^T\}$, кососимметрические ---
$\{A\in M_n(F)\mid A=-A^T\}$.


При $\dim U=n<\infty$ второе сопряжённое пространство $U^{**}$ также
имеет размерность $n$. Существует канонический изоморфизм
\[
  \iota\colon U\to U^{**},
  \qquad
  v\mapsto f_v,
  \qquad
  f_v(\omega)=\omega(v).
\]
Он не зависит от выбора базиса: линейность следует из
\[
  f_{v+w}=f_v+f_w,
  \qquad f_{\lambda v}=\lambda f_v,
\]
а в конечномерном случае размерности совпадают.

\begin{proof}
Если $v\ne0$, дополним его до базиса $U$ и возьмём координатный
функционал $\omega$, равный $1$ на $v$ и $0$ на остальных базисных
векторах. Тогда $f_v(\omega)=\omega(v)\ne0$, поэтому $\iota$ инъективно.
В конечномерном случае
$\dim U=\dim U^{**}$, следовательно, инъективное отображение $\iota$
сюръективно и является изоморфизмом.
\end{proof}

\begin{definition}
Для подпространства $U\leq V$ аннулятором называется
\[
  U^0=\{\omega\in V^*\mid \omega(u)=0\text{ для всех }u\in U\}
  \leq V^*.
\]
\end{definition}

Если $\dim V=n$ и $\dim U=k$, то
\[
  \dim U^0=n-k=\operatorname{codim}U,
\]
и $U^0\cong(V/U)^*$.

\begin{proof}
Выберем базис $e_1,\ldots,e_k$ пространства $U$ и дополним его до
базиса $e_1,\ldots,e_n$ пространства $V$. Тогда
$e^{k+1},\ldots,e^n$ образуют базис $U^0$, откуда
$\dim U^0=n-k$. Отображение
\[
  (V/U)^*\to U^0,\qquad
  \lambda\mapsto\lambda\circ\pi,
\]
где $\pi\colon V\to V/U$ --- каноническая проекция, является
изоморфизмом.
\end{proof}


Из формулы размерности следует
\[
  \dim (U^0)^0=\dim U,
\]
а поскольку $U\subseteq(U^0)^0$, в конечномерном случае
\[
  (U^0)^0=U.
\]

\begin{theorem}[Критерий базиса в $V^*$]
Функционалы $\omega_1,\ldots,\omega_n$ образуют базис $V^*$ тогда и
только тогда, когда
\[
  \bigcap_{i=1}^{n}\Ker\omega_i=\{0\}.
\]
Эквивалентно, их линейная оболочка имеет нулевой аннулятор.
\end{theorem}

\begin{proof}
Рассмотрим отображение
\[
  \Phi\colon V\to F^n,\qquad
  v\mapsto(\omega_1(v),\ldots,\omega_n(v)).
\]
Его ядро равно $\bigcap_i\Ker\omega_i$. Если это пересечение нулевое,
то $\Phi$ инъективно. Так как $\dim V=\dim F^n=n$, отображение
изоморфно, а его координатные функционалы $\omega_i$ образуют базис
$V^*$. Обратно, если $\omega_i$ образуют базис и все они обращаются в
нуль на $v$, то на $v$ обращается в нуль любой функционал. Координатный
функционал относительно базиса, содержащего ненулевой $v$, показывает,
что это возможно только при $v=0$.
\end{proof}

\section{Сопряжённый гомоморфизм}

Пусть $A\colon U\to V$. Сопряжённое отображение
\[
  A^*\colon V^*\to U^*,
  \qquad
  A^*(\omega)=\omega\circ A
\]
задаётся равенством
\[
  (A^*\omega)(v)=\omega(Av).
\]
Оно обладает свойствами
\[
  (A+B)^*=A^*+B^*,
  \qquad (\lambda A)^*=\lambda A^*,
  \qquad (AB)^*=B^*A^*.
\]
Матрица сопряжённого отображения в сопряжённых базисах является
транспонированной матрицей исходного отображения.

\begin{visualreconstruction}
Определение удобно представить цепочкой стрелок
$A\colon U\to V$, $A^*\colon V^*\to U^*$ и соответствием
$\omega\mapsto\omega\circ A$:
\[
  (A^*(\omega))(v)=(A^*\omega,v)=(\omega,Av)
  =\omega(A(v)),
  \qquad v\in U,\ \omega\in V^*.
\]
\end{visualreconstruction}


Имеют место соотношения
\[
  \Ker A^*=(\Im A)^0,
  \qquad
  \Im A^*=(\Ker A)^0.
\]

\begin{proof}
Функционал $\omega\in V^*$ лежит в $\Ker A^*$ тогда и только тогда,
когда $\omega(Au)=0$ для всех $u\in U$, то есть когда он обращается в
нуль на $\Im A$. Поэтому $\Ker A^*=(\Im A)^0$.

Для второго равенства включение
$\Im A^*\subseteq(\Ker A)^0$ следует из
$(A^*\omega)(u)=\omega(Au)=0$ при $u\in\Ker A$. Размерности этих
пространств совпадают:
\[
  \dim\Im A^*=\rk A^*=\rk A
  =\dim U-\dim\Ker A=\dim(\Ker A)^0.
\]
Следовательно, включение является равенством.
\end{proof}
В частности,
\[
  \dim\Im A
  =\dim U-\dim\Ker A
  =\dim(\Ker A)^0
  =\dim\Im A^*,
\]
что ещё раз доказывает равенство строчного и столбцового рангов матрицы.


\chapter{Тензоры}

Пусть $V$ --- конечномерное векторное пространство, а $V^*$ ---
сопряжённое.

\begin{definition}
Тензор типа $(p,q)$ --- полилинейное отображение
\[
  T\colon V^p\times(V^*)^q\to F.
\]
Число $p+q$ называется полной валентностью (рангом) тензора; $p$ ---
ковариантная, $q$ --- контравариантная валентность.
Множество таких тензоров обозначается $T_p^q(V)$.
\end{definition}

\begin{explanation}
Небольшая схема $v=2e_1$, $e_1=\tfrac12v$ иллюстрирует главный мотив
раздела: при замене базиса меняются координаты аргументов и компонент,
но сам тензор как полилинейное отображение остаётся тем же.
\end{explanation}

Частные случаи:
\begin{enumerate}[label=\arabic*)]
  \item тензоры типа $(0,0)$ --- скаляры $F$;
  \item типа $(1,0)$ --- ковекторы $V^*$;
  \item типа $(0,1)$ --- векторы $(V^*)^*\cong V$;
  \item типа $(2,0)$ --- билинейные функции на $V$;
  \item типа $(1,1)$ --- линейные операторы $V\to V$;
  \item типа $(n,0)$ --- в частности, определитель квадратной матрицы.
\end{enumerate}

Для $T,S\in T_p^q(V)$ определяются сумма и умножение на скаляр
поточечно. Тензорное произведение
$T\otimes S\in T_{p+k}^{q+\ell}(V)$ для
$S\in T_k^{\ell}(V)$ задаётся формулой
\[
\begin{aligned}
&(T\otimes S)(v_1,\ldots,v_{p+k},
\omega_1,\ldots,\omega_{q+\ell})\\
&\qquad=T(v_1,\ldots,v_p,\omega_1,\ldots,\omega_q)
S(v_{p+1},\ldots,v_{p+k},
\omega_{q+1},\ldots,\omega_{q+\ell}).
\end{aligned}
\]


Тензорное произведение ассоциативно, вообще говоря некоммутативно,
дистрибутивно относительно сложения и согласовано с умножением на скаляры:
\[
  T\otimes(S\otimes R)=(T\otimes S)\otimes R,
\]
\[
  T\otimes(S+R)=T\otimes S+T\otimes R,
\]
\[
  \lambda(T\otimes S)=(\lambda T)\otimes S=T\otimes(\lambda S).
\]

\begin{theorem}[Тензорный базис]
Если $e_1,\ldots,e_n$ --- базис $V$, а $e^1,\ldots,e^n$ --- сопряжённый
базис $V^*$, то тензоры
\[
  e^{i_1}\otimes\cdots\otimes e^{i_p}
  \otimes e_{j_1}\otimes\cdots\otimes e_{j_q}
\]
образуют базис $T_p^q(V)$.
\end{theorem}

\begin{proof}
Полилинейность показывает, что тензор однозначно определяется значениями
на наборах базисных векторов и ковекторов. Поэтому любой тензор
раскладывается по указанной системе с коэффициентами
\[
  T_{i_1\ldots i_p}^{j_1\ldots j_q}
  =T(e_{i_1},\ldots,e_{i_p},
      e^{j_1},\ldots,e^{j_q}).
\]
Линейная независимость проверяется подстановкой тех же базисных
аргументов: равенство $e^i(e_j)=\delta^i_j$ выделяет ровно один
коэффициент линейной комбинации. Следовательно, система одновременно
порождает пространство и линейно независима.
\end{proof}

Отсюда
\[
  \dim T_p^q(V)=n^{p+q}.
\]

\begin{example}
Пусть $\dim V=2$. При
базисах $e_1,e_2$ и $e^1,e^2$ четыре тензора
\[
  e^1\otimes e_1,\quad e^1\otimes e_2,\quad
  e^2\otimes e_1,\quad e^2\otimes e_2
\]
подставляются в предполагаемую линейную комбинацию с коэффициентами
$\lambda^i_j$. Подстановка пары $(e_k,e^\ell)$ выделяет коэффициент по
правилу
\[
  (e^i\otimes e_j)(e_k,e^\ell)
  =e^i(e_k)e^\ell(e_j)=\delta^i_k\delta^\ell_j.
\]
Отсюда каждый $\lambda^i_j=0$.
\end{example}


Компоненты тензора в выбранном базисе находятся подстановкой базисных
векторов и ковекторов. Например, для $T\in T_1^1(V)$
\[
  T=\sum_{i,j}T^i{}_j\,e^j\otimes e_i,
\]
и числа $T^i{}_j$ образуют матрицу. Для $T\in T_2^0(V)$ компоненты
$T_{ij}$ также образуют квадратную матрицу.


У тензора второго ранга встречаются четыре вида расположения индексов:
\[
  T^{ij}\in T_0^2(V),\qquad
  T^i{}_j\in T_1^1(V),\qquad
  T_i{}^j\in T_1^1(V),\qquad
  T_{ij}\in T_2^0(V).
\]
Для тензора третьего ранга соответственно возможны компоненты
\[
  T^{ijk},\qquad T^{ij}{}_k,\qquad
  T^i{}_{jk},\qquad T_{ijk},
\]
то есть тензоры типов $(0,3)$, $(1,2)$, $(2,1)$ и $(3,0)$.
Аналогично записываются компоненты тензоров четвёртого и более высоких
рангов: число нижних индексов равно первой валентности, число верхних ---
второй.

В матричном представлении сложению тензоров соответствует сложение
матриц, а умножению тензора на скаляр --- умножение матрицы на этот
скаляр.

Для $\dim V=3$ компоненты тензоров типов $(1,1)$ и $(2,0)$ удобно
записывать матрицами:
\[
  T\in T_1^1(V)\ \longleftrightarrow
  \begin{pmatrix}
    T^1{}_1&T^1{}_2&T^1{}_3\\
    T^2{}_1&T^2{}_2&T^2{}_3\\
    T^3{}_1&T^3{}_2&T^3{}_3
  \end{pmatrix},
  \qquad
  T\in T_2^0(V)\ \longleftrightarrow
  \begin{pmatrix}
    T_{11}&T_{12}&T_{13}\\
    T_{21}&T_{22}&T_{23}\\
    T_{31}&T_{32}&T_{33}
  \end{pmatrix}.
\]

Рассмотрим замену базиса. Пусть
\[
  \widetilde e_j=\sum_i e_i C^i{}_j,
\]
то есть столбцы матрицы $C$ состоят из координат нового базиса в старом.
Тогда для любого $x\in V$
\[
  [x]_e=C[x]_{\widetilde e},
  \qquad
  [x]_{\widetilde e}=C^{-1}[x]_e,
\]
а сопряжённый базис преобразуется по закону
\[
  \widetilde e^{\,i}=\sum_k(C^{-1})^i{}_k e^k.
\]


Пусть $T\in T_1^1(V)$ и $A=(T^i{}_j)$ --- матрица его компонент в
базисе $e$. В новом базисе
\[
  \widetilde T^{\,i}{}_{j}
  =\sum_{k,\ell}(C^{-1})^i{}_k T^k{}_{\ell}C^{\ell}{}_j,
\]
или в матричной записи
\[
  \widetilde A=C^{-1}AC.
\]

\begin{definition}
Формула, выражающая компоненты тензора в новом базисе через его
компоненты в старом, называется \emph{тензорным законом преобразования
координат}. В общем случае каждый нижний индекс преобразуется матрицей
$C$, а каждый верхний --- матрицей $C^{-1}$.
\end{definition}


\begin{definition}
Тензор называется \emph{инвариантным}, если он имеет одинаковые
компоненты во всех тензорных базисах.
\end{definition}

Существует естественный изоморфизм
\[
  \operatorname{End}(V)\cong T_1^1(V).
\]
Линейному оператору $A\colon V\to V$ соответствует тензор
\[
  T(v,\omega)=\omega(Av),
  \qquad v\in V,\quad \omega\in V^*.
\]
Действительно, если
\[
  Ae_j=\sum_i e_i a^i{}_j,
\]
то
\[
  T(e_j,e^i)=e^i(Ae_j)=a^i{}_j.
\]
Следовательно,
\[
  T_1^1(V)\cong M_n(F)\cong\operatorname{End}(V).
\]

Соответствия удобно свести в двунаправленную схему:
\[
  T_1^1(V)
  \underset{T\longmapsto(a^i{}_j)}{\stackrel{\sim}{\longleftrightarrow}}
  M_n(F)
  \underset{A\longmapsto(a^i{}_j)}{\stackrel{\sim}{\longleftrightarrow}}
  \operatorname{End}(V).
\]

\section{Симметрические и кососимметрические тензоры}

Тензор $T\in T_p^0(V)$ называется \emph{симметрическим}, если для любой
перестановки $\sigma\in S_p$
\[
  T(x_{\sigma(1)},\ldots,x_{\sigma(p)})
  =T(x_1,\ldots,x_p).
\]
Он называется \emph{кососимметрическим} (альтернированным), если
\[
  T(x_{\sigma(1)},\ldots,x_{\sigma(p)})
  =\sgn(\sigma)T(x_1,\ldots,x_p).
\]


Пусть $\operatorname{char}F=0$. Для $T\in T_p^0(V)$ определяются
симметризация и альтернирование:
\[
  \operatorname{Sym}T(x_1,\ldots,x_p)
  =\frac1{p!}\sum_{\sigma\in S_p}
  T(x_{\sigma(1)},\ldots,x_{\sigma(p)}),
\]
\[
  \operatorname{Alt}T(x_1,\ldots,x_p)
  =\frac1{p!}\sum_{\sigma\in S_p}\sgn(\sigma)
  T(x_{\sigma(1)},\ldots,x_{\sigma(p)}).
\]
Операторы $\operatorname{Sym}$ и $\operatorname{Alt}$ линейны и
идемпотентны:
\[
  \operatorname{Sym}^2=\operatorname{Sym},
  \qquad
  \operatorname{Alt}^2=\operatorname{Alt}.
\]

\begin{proof}
Линейность следует из линейности суммы. При повторной симметризации
каждая перестановка аргументов возникает ровно $p!$ раз, поэтому второй
множитель $1/p!$ сокращает число повторений и
$\operatorname{Sym}^2=\operatorname{Sym}$. Для альтернирования
используется дополнительно
$\sgn(\sigma\tau)=\sgn(\sigma)\sgn(\tau)$; после перегруппировки суммы
каждое слагаемое снова появляется $p!$ раз с требуемым знаком.
\end{proof}
Если $T$ симметричен, то
$\operatorname{Sym}T=T$ и $\operatorname{Alt}T=0$; если $T$
кососимметричен, то
$\operatorname{Alt}T=T$ и $\operatorname{Sym}T=0$. Поэтому образы этих
операторов состоят соответственно из симметрических и кососимметрических
тензоров.

\begin{editorialnote}
Равенства $\operatorname{Alt}T=0$ для симметрического $T$ и
$\operatorname{Sym}T=0$ для кососимметрического $T$ требуют $p\ge2$
(и делимости на $p!$, обеспеченной условием
$\operatorname{char}F=0$). При $p=1$ оба оператора совпадают с
тождественным.
\end{editorialnote}

Симметрическое произведение симметрических тензоров $T$ и $S$ задаётся
формулой
\[
  T\odot S=\operatorname{Sym}(T\otimes S).
\]
Симметризованные произведения элементов сопряжённого базиса образуют
базис пространства симметрических тензоров степени $p$. Если
$n=\dim V$, то
\[
  \dim S^p(V^*)=\binom{n+p-1}{p}.
\]

\begin{explanation}
Биномиальный коэффициент $\binom{n+p-1}{p}$ имеет комбинаторный смысл:
это число способов распределить $p$ неразличимых множителей по $n$
базисным ковекторам, то есть число мономов степени $p$ от $n$
переменных.
\end{explanation}


\chapter*{Примеры групп}
\addcontentsline{toc}{chapter}{Примеры групп}

\begin{enumerate}
  \item $V_4$ --- четверная группа Клейна.
  \item $A_n$ --- знакопеременная группа, то есть группа чётных
  перестановок.
  \item Группа углов с операцией сложения углов.
  \item Аддитивная группа векторов плоскости.
  \item Группа параллельных переносов плоскости.
  \item $\Z_n$ --- группа классов вычетов по модулю $n$.
  \item $R_n$ и $C_n$ --- группы поворотов правильного $n$-угольника.
  Обозначение второго семейства восстановлено как
  \uncertain{$C_n$; графема допускает иное чтение}.
  \item $D_n$ --- диэдральная группа всех симметрий правильного
  $n$-угольника; $|R_n|=n$, $|D_n|=2n$.
  \item $S_n$ --- группа всех перестановок $n$ элементов.
  \item $\mu_n$ --- группа корней $n$-й степени из единицы:
  \[
    \mu_n=\left\{e^{2\pi i k/n}\mid k=0,\ldots,n-1\right\}
    \cong\Z_n.
  \]
  \item Общая линейная группа
  \[
    GL_n(F)=\{A\in M_n(F)\mid\det A\ne0\}.
  \]
  \item Специальная линейная группа
  \[
    SL_n(F)=\{A\in M_n(F)\mid\det A=1\}
    \trianglelefteq GL_n(F).
  \]
  \item $T_n(F)$ --- группа невырожденных верхнетреугольных матриц.
  \item $UT_n(F)$ --- группа верхних унитреугольных матриц.
  \item $D_n(F)$ --- группа невырожденных диагональных матриц.
  \item $SO_n(\R)$ --- специальная ортогональная группа:
  \[
    SO_n(\R)=\{A\in M_n(\R)\mid A^{\mathsf T}A=E,\ \det A=1\}.
  \]
\end{enumerate}


\appendix
\chapter{Редакторские пометки}

Большинство повреждённых промежуточных выкладок восстановлено по
следующим за ними утверждениям и стандартным доказательствам. В основном
тексте такие реконструкции не смешаны с авторскими формулировками:
добавленные рассуждения находятся в окружениях «Доказательство» и
«Пояснение».

\begin{longtable}{@{}p{0.28\textwidth}p{0.64\textwidth}@{}}
\toprule
Место & Принятое редакторское решение \\
\midrule
\endfirsthead
\toprule
Место & Принятое редакторское решение \\
\midrule
\endhead
Фундаментальная транспозиция &
Повреждённый подсчёт числа соседних транспозиций восстановлен как
$2(j-i)-1$. Это единственная формула, согласующаяся с изображённым
разложением и нечётностью транспозиции. \\
Критерий зависимости для модулей &
Исходная формулировка сохранена, но добавлена оговорка: над общим кольцом
она требует обратимости одного из коэффициентов. Безусловная версия
доказана для векторных пространств над полем. \\
Простой и неприводимый элементы &
Отмечено различие стандартных определений. Их совпадение используется
только после доказательства факториальности евклидова кольца. \\
Подстановка в многочлен &
Уточнено условие коммутативности кольца либо центральности подставляемого
элемента. \\
Симметризация и альтернирование &
Уточнено, что взаимное зануление операторов для симметрических и
кососимметрических тензоров предполагает $p\ge2$. \\
Обозначение группы поворотов &
Символ восстановлен как $C_n$, но чтение остаётся не вполне надёжным.
Обычно $C_n$ обозначает циклическую группу поворотов, тогда как $R_n$ не
является универсально принятым обозначением. \\
\bottomrule
\end{longtable}

\begin{editorialnote}
Полностью зачёркнутые слова, случайные пробы пера и неподписанные
декоративные штрихи в учебный текст не включались. Там, где они могли
скрывать математический переход, вместо догадки приведено самостоятельное
доказательство соответствующего результата.
\end{editorialnote}

\end{document}